\IfFileExists{stacks-project.cls}{%
\documentclass{stacks-project}
}{%
\documentclass{amsart}
}
\usepackage{amssymb}

% For dealing with references we use the comment environment
\usepackage{verbatim}
\newenvironment{reference}{\comment}{\endcomment}
%\newenvironment{reference}{}{}
\newenvironment{slogan}{\comment}{\endcomment}
\newenvironment{history}{\comment}{\endcomment}

% For commutative diagrams we use Xy-pic
\usepackage[all]{xy}

% We use 2cell for 2-commutative diagrams.
\xyoption{2cell}
\UseAllTwocells

% We use multicol for the list of chapters between chapters
\usepackage{multicol}

% This is generally recommended for better output
\usepackage{lmodern}
\usepackage[T1]{fontenc}
\usepackage{textalpha}

% For cross-file-references
% We prefer the snapshot tmp/xref/#2.aux,
% which latexmkrc replaces atomically, so
% that documents compiled simultaneously
% never read a half-written .aux file.
% The last argument keeps links to #2.pdf.
\usepackage{xr-hyper}
\newcommand{\maybeexternaldocument}[2]{%
  \IfFileExists{#2.aux}{%
    \externaldocument[#1]{#2}%
  }{%
    \PackageWarning{stack}{Missing #2.aux; external references with prefix #1 unavailable}%
  }%
}
\let\externaldocumentorig\externaldocument
\renewcommand{\externaldocument}[2][]{%
  \IfFileExists{tmp/xref/#2.aux}{%
    \externaldocumentorig[#1]{tmp/xref/#2}[#2.pdf]%
  }{%
  \IfFileExists{#2.aux}{%
    \externaldocumentorig[#1]{#2}%
  }{%
    \PackageWarning{stack}{Missing #2.aux; external references with prefix #1 unavailable}%
  }%
  }%
}

% Package for hypertext links:
\usepackage{hyperref}

% For any local file, say "hello.tex" you want to link to please
% use \externaldocument[hello-]{hello}
\externaldocument[introduction-]{introduction}
\externaldocument[conventions-]{conventions}
\externaldocument[sets-]{sets}
\externaldocument[categories-]{categories}
\externaldocument[topology-]{topology}
\externaldocument[sheaves-]{sheaves}
\externaldocument[sites-]{sites}
\externaldocument[stacks-]{stacks}
\externaldocument[fields-]{fields}
\externaldocument[algebra-]{algebra}
\externaldocument[brauer-]{brauer}
\externaldocument[homology-]{homology}
\externaldocument[derived-]{derived}
\externaldocument[simplicial-]{simplicial}
\externaldocument[more-algebra-]{more-algebra}
\externaldocument[smoothing-]{smoothing}
\externaldocument[modules-]{modules}
\externaldocument[sites-modules-]{sites-modules}
\externaldocument[injectives-]{injectives}
\externaldocument[cohomology-]{cohomology}
\externaldocument[sites-cohomology-]{sites-cohomology}
\externaldocument[dga-]{dga}
\externaldocument[dpa-]{dpa}
\externaldocument[sdga-]{sdga}
\externaldocument[hypercovering-]{hypercovering}
\externaldocument[schemes-]{schemes}
\externaldocument[constructions-]{constructions}
\externaldocument[properties-]{properties}
\externaldocument[morphisms-]{morphisms}
\externaldocument[coherent-]{coherent}
\externaldocument[divisors-]{divisors}
\externaldocument[limits-]{limits}
\externaldocument[varieties-]{varieties}
\externaldocument[topologies-]{topologies}
\externaldocument[descent-]{descent}
\externaldocument[perfect-]{perfect}
\externaldocument[more-morphisms-]{more-morphisms}
\externaldocument[flat-]{flat}
\externaldocument[groupoids-]{groupoids}
\externaldocument[more-groupoids-]{more-groupoids}
\externaldocument[etale-]{etale}
\externaldocument[chow-]{chow}
\externaldocument[intersection-]{intersection}
\externaldocument[pic-]{pic}
\externaldocument[weil-]{weil}
\externaldocument[adequate-]{adequate}
\externaldocument[dualizing-]{dualizing}
\externaldocument[duality-]{duality}
\externaldocument[discriminant-]{discriminant}
\externaldocument[derham-]{derham}
\externaldocument[local-cohomology-]{local-cohomology}
\externaldocument[algebraization-]{algebraization}
\externaldocument[curves-]{curves}
\externaldocument[resolve-]{resolve}
\externaldocument[models-]{models}
\externaldocument[functors-]{functors}
\externaldocument[equiv-]{equiv}
\externaldocument[pione-]{pione}
\externaldocument[etale-cohomology-]{etale-cohomology}
\externaldocument[proetale-]{proetale}
\externaldocument[relative-cycles-]{relative-cycles}
\externaldocument[more-etale-]{more-etale}
\externaldocument[trace-]{trace}
\externaldocument[crystalline-]{crystalline}
\externaldocument[spaces-]{spaces}
\externaldocument[spaces-properties-]{spaces-properties}
\externaldocument[spaces-morphisms-]{spaces-morphisms}
\externaldocument[decent-spaces-]{decent-spaces}
\externaldocument[spaces-cohomology-]{spaces-cohomology}
\externaldocument[spaces-limits-]{spaces-limits}
\externaldocument[spaces-divisors-]{spaces-divisors}
\externaldocument[spaces-over-fields-]{spaces-over-fields}
\externaldocument[spaces-topologies-]{spaces-topologies}
\externaldocument[spaces-descent-]{spaces-descent}
\externaldocument[spaces-perfect-]{spaces-perfect}
\externaldocument[spaces-more-morphisms-]{spaces-more-morphisms}
\externaldocument[spaces-flat-]{spaces-flat}
\externaldocument[spaces-groupoids-]{spaces-groupoids}
\externaldocument[spaces-more-groupoids-]{spaces-more-groupoids}
\externaldocument[bootstrap-]{bootstrap}
\externaldocument[spaces-pushouts-]{spaces-pushouts}
\externaldocument[spaces-chow-]{spaces-chow}
\externaldocument[groupoids-quotients-]{groupoids-quotients}
\externaldocument[spaces-more-cohomology-]{spaces-more-cohomology}
\externaldocument[spaces-simplicial-]{spaces-simplicial}
\externaldocument[spaces-duality-]{spaces-duality}
\externaldocument[formal-spaces-]{formal-spaces}
\externaldocument[restricted-]{restricted}
\externaldocument[spaces-resolve-]{spaces-resolve}
\externaldocument[formal-defos-]{formal-defos}
\externaldocument[defos-]{defos}
\externaldocument[cotangent-]{cotangent}
\externaldocument[examples-defos-]{examples-defos}
\externaldocument[algebraic-]{algebraic}
\externaldocument[examples-stacks-]{examples-stacks}
\externaldocument[stacks-sheaves-]{stacks-sheaves}
\externaldocument[criteria-]{criteria}
\externaldocument[artin-]{artin}
\externaldocument[quot-]{quot}
\externaldocument[stacks-properties-]{stacks-properties}
\externaldocument[stacks-morphisms-]{stacks-morphisms}
\externaldocument[stacks-limits-]{stacks-limits}
\externaldocument[stacks-cohomology-]{stacks-cohomology}
\externaldocument[stacks-perfect-]{stacks-perfect}
\externaldocument[stacks-introduction-]{stacks-introduction}
\externaldocument[stacks-more-morphisms-]{stacks-more-morphisms}
\externaldocument[stacks-geometry-]{stacks-geometry}
\externaldocument[moduli-]{moduli}
\externaldocument[moduli-curves-]{moduli-curves}
\externaldocument[examples-]{examples}
\externaldocument[exercises-]{exercises}
\externaldocument[guide-]{guide}
\externaldocument[desirables-]{desirables}
\externaldocument[coding-]{coding}
\externaldocument[obsolete-]{obsolete}
\externaldocument[fdl-]{fdl}
\externaldocument[index-]{index}

% Heap Project documents
\externaldocument[linear-algebra-]{linear-algebra}
\externaldocument[groups-]{groups}
\externaldocument[complex-analysis-]{complex-analysis}
\externaldocument[functional-analysis-]{functional-analysis}
\externaldocument[categories-]{categories}
\externaldocument[commutative-algebra-]{commutative-algebra}
\externaldocument[ringed-spaces-]{ringed-spaces}
\externaldocument[homological-algebra-]{homological-algebra}
\externaldocument[lie-algebras-]{lie-algebras}
\externaldocument[representation-theory-]{representation-theory}
\externaldocument[smooth-manifolds-]{smooth-manifolds}
\externaldocument[differential-topology-]{differential-topology}
\externaldocument[algebraic-topology-]{algebraic-topology}
\externaldocument[differential-geometry-]{differential-geometry}
\externaldocument[algebraic-geometry-]{algebraic-geometry}
\externaldocument[symplectic-geometry-]{symplectic-geometry}
\externaldocument[complex-geometry-]{complex-geometry}
\externaldocument[hodge-theory-]{hodge-theory}
\maybeexternaldocument{deformations-}{deformations}
\externaldocument[vertex-algebras-]{vertex-algebras}
\externaldocument[springer-]{springer}
\externaldocument[infty-categories-]{infty-categories}
\externaldocument[seminars-differential-geometry-]{%
seminars-differential-geometry}
\externaldocument[seminars-complex-geometry-]{%
seminars-complex-geometry}
\externaldocument[seminars-algebraic-geometry-]{%
seminars-algebraic-geometry}
\externaldocument[seminars-others-]{%
seminars-others}
\externaldocument[%
geometric-prequantization-]{%
geometric-prequantization}
\externaldocument[surfaces-]{surfaces}
\externaldocument[birational-maps-commutative-algebra-]{birational-maps-commutative-algebra}
\externaldocument[cremona-transformations-]{cremona-transformations}
\externaldocument[derived-categories-of-sheaves-]{derived-categories-of-sheaves}
\externaldocument[geometric-invariant-theory-]{geometric-invariant-theory}
\externaldocument[geometric-stability-conditions-and-group-actions-]{geometric-stability-conditions-and-group-actions}
\externaldocument[moduli-spaces-of-sheaves-]{moduli-spaces-of-sheaves}
\externaldocument[bridgeland-stability-general-theory-]{bridgeland-stability-general-theory}
\externaldocument[hilbert-schemes-]{hilbert-schemes}
\externaldocument[noncommutative-abelian-surfaces-and-kummer-type-hyperkahler-varieties-]{noncommutative-abelian-surfaces-and-kummer-type-hyperkahler-varieties}
\externaldocument[mumford-tate-groups-in-hodge-theory-]{mumford-tate-groups-in-hodge-theory}
\externaldocument[hodge-birational-atoms-2026-]{hodge-birational-atoms-2026}
\externaldocument[xxii-egd-2026-]{%
xxii-egd-2026}
\externaldocument[pde-]{pde}
\externaldocument[probability-]{probability}

\externaldocument[linguistics-]{linguistics}
\externaldocument[genealogia-familiar-]{%
genealogia-familiar}

\externaldocument[%
16th-alga-meeting-2026-]{%
16th-alga-meeting-2026}

\externaldocument[%
iv-jornadas-lefschetz-2026-]{%
iv-jornadas-lefschetz-2026}

\externaldocument[reading-the-masters-]{%
reading-the-masters}

\externaldocument[real-analysis-]{%
real-analysis}

\externaldocument[optimal-transport-]{%
optimal-transport}

\externaldocument[history-]{history}

\externaldocument[seiberg-witten-]{seiberg-witten}

% Theorem environments.
%
\theoremstyle{plain}
\newtheorem{theorem}[subsection]{Theorem}
\newtheorem{proposition}[subsection]{Proposition}
\newtheorem{lemma}[subsection]{Lemma}

\theoremstyle{definition}
\newtheorem{definition}[subsection]{Definition}
\newtheorem{example}[subsection]{Example}
\newtheorem{exercise}[subsection]{Exercise}
\newtheorem{situation}[subsection]{Situation}

\theoremstyle{remark}
\newtheorem{remark}[subsection]{Remark}
\newtheorem{remarks}[subsection]{Remarks}

\numberwithin{equation}{subsection}

% Macros
%
\def\lim{\mathop{\mathrm{lim}}\nolimits}
\def\colim{\mathop{\mathrm{colim}}\nolimits}
\def\Spec{\mathop{\mathrm{Spec}}}
\def\Hom{\mathop{\mathrm{Hom}}\nolimits}
\def\Ext{\mathop{\mathrm{Ext}}\nolimits}
\def\SheafHom{\mathop{\mathcal{H}\!\mathit{om}}\nolimits}
\def\SheafExt{\mathop{\mathcal{E}\!\mathit{xt}}\nolimits}
\def\Sch{\mathit{Sch}}
\def\Mor{\mathop{\mathrm{Mor}}\nolimits}
\def\Ob{\mathop{\mathrm{Ob}}\nolimits}
\def\Sh{\mathop{\mathit{Sh}}\nolimits}
\def\NL{\mathop{N\!L}\nolimits}
\def\CH{\mathop{\mathrm{CH}}\nolimits}
\def\proetale{{pro\text{-}\acute{e}tale}}
\def\etale{{\acute{e}tale}}
\def\QCoh{\mathit{QCoh}}
\def\Ker{\mathop{\mathrm{Ker}}}
\def\Im{\mathop{\mathrm{Im}}}
\def\Coker{\mathop{\mathrm{Coker}}}
\def\Coim{\mathop{\mathrm{Coim}}}

% Boxtimes
%
\DeclareMathSymbol{\boxtimes}{\mathbin}{AMSa}{"02}

%
% Macros for moduli stacks/spaces
%
\def\QCohstack{\mathcal{QC}\!\mathit{oh}}
\def\Cohstack{\mathcal{C}\!\mathit{oh}}
\def\Spacesstack{\mathcal{S}\!\mathit{paces}}
\def\Quotfunctor{\mathrm{Quot}}
\def\Hilbfunctor{\mathrm{Hilb}}
\def\Curvesstack{\mathcal{C}\!\mathit{urves}}
\def\Polarizedstack{\mathcal{P}\!\mathit{olarized}}
\def\Complexesstack{\mathcal{C}\!\mathit{omplexes}}
% \Pic is the operator that assigns to X its picard group, usage \Pic(X)
% \Picardstack_{X/B} denotes the Picard stack of X over B
% \Picardfunctor_{X/B} denotes the Picard functor of X over B
\def\Pic{\mathop{\mathrm{Pic}}\nolimits}
\def\Picardstack{\mathcal{P}\!\mathit{ic}}
\def\Picardfunctor{\mathrm{Pic}}
\def\Deformationcategory{\mathcal{D}\!\mathit{ef}}


\begin{document}

\title{Algebraic Topology}
\maketitle

\phantomsection
\label{section-phantom}

\tableofcontents


\medskip
\noindent
{\bf Organization of these notes.}
These notes combine material from several courses and sets of
lecture notes.  The first block develops homotopy theory of spaces
(coverings, homotopy groups, fibrations, cofibrations, and model
structures).  The second block develops the standard structural
results (Blakers--Massey, Freudenthal, Whitehead and Postnikov
towers, homotopy fibers and cofibers, and representability).
The final block introduces spectral sequences, especially the
Serre spectral sequence and its applications.  Some subsections
are still lecture-note fragments; these are left visibly marked
rather than silently completed.


\section{Covering spaces}
\label{section-covering-spaces}

\begin{proposition}[Unique lifting property]
\label{proposition-unique-lifting-property}
\cite{hat}, Proposition 1.34. Given a
covering space $p:\tilde{X} \to X$ and a
map $f:Y \to X$, if two lifts
$\tilde{f}_1,\tilde{f}_2:Y \to \tilde{X}$ of
$f$
agree at one point of $Y$ and $Y$ is
connected, then $\tilde{f}_1$ and
$\tilde{f}_2$ agree on all of $Y$.
\end{proposition}

\begin{proof}
We aim to show that the set of points where
$\tilde{f}_1$ and $\tilde{f}_2$
agree is both open and closed in $Y$. Let $y
\in Y$ and consider an evenly
covered neighbourhood $U$ of $f(y)$, so that
$p^{-1}(U)$ is a disjoint union of
sets, each of which is homeomorphic to $U$
via $p$. Let $\tilde{U}_1$ and
$\tilde{U}_2$ be the sheets containing
$\tilde{f}_1(y)$ and $\tilde{f}_2(y)$,
respectively. By continuity of $\tilde{f}_1$
and $\tilde{f}_2$, the intersection
$\tilde{f}^{-1}_1(\tilde{U}_1) \cap
\tilde{f}_2^{-1}(\tilde{U}_2):=N$ is
 an open neighbourhood of $y$.

Since  $\tilde{U}_1$ and $\tilde{U}_2$ are
mapped homeomorphically to the set
$U \subset X$, we can only have that either
$\tilde{U}_1$ coincides
with $\tilde{U}_2$, or they are disjoint. If
they are disjoint, then the
neighbourhood $N$ shows that the set of
points where $\tilde{f}_1$ and
$\tilde{f}_2$ agree is closed; that's because
in this case they don't agree at a
point, so that we are on the complement of
the set where they agree, and we have
shown it's an open set (every point has a
neighbourhood contained in the set).

If, on the other hand, $\tilde{U}_1$ and
$\tilde{U}_2$ are the same, which
happens when $\tilde{f}_1(y)=\tilde{f}_2(y)$,
then $\tilde{f}_1$ and
$\tilde{f}_2$ agree throughout $N$, so that
we have shown that the set where
$\tilde{f}_1$ and $\tilde{f}_2$ coincide is
open.
\end{proof}

\begin{definition}[Properly discontinuous
action]
\label{definition-prope-disco-actio}
See \cite{hat} p. 72. The action of a group
$G$ on a space $Y$ is called
{\it properly discontinuous} or {\it covering
space action} if each $y \in Y$
has a neighbourhood $U$ such that $g_1(U)
\cap g_2(U)\neq \emptyset$ for any
$g_1,g_2 \in G$ implies $g_1=g_2$.
\end{definition}

\begin{proposition}
\label{proposition-deck-trans-act-prope}
The action of deck transformations on a
covering space is a properly
discontinuous action.
\end{proposition}

\section{Definitions in homotopy theory}
\label{section-defin-homot-theor}

This is taken from \cite{TomDieck}.

\begin{itemize}
\item Let $X$ and $Y$ be topological spaces
and
 $f,g:X\to Y$ continuous maps.
An {\it homotopy} from $f$ to $g$ is a
continuous map
$$
H:X\times[0,1]\to Y,\qquad(x,t)\mapsto
H(x,t)=H_t(x)
$$
\noindent
such that $f(x)=H(x,0)$ and $g(x)=H(x,1)$ for
all $x\in X$.
We denote this situation by $f\simeq g$.
The homotopy relation $\simeq$ is an
equivalence relation
on the set of continuous maps $X\to Y$.
A homotopy of maps $H_t:X\to Y$ is called
{\it relative to $A\subset X$} if $H_t|_A$ is
constant.


\item Topological spaces and homotopy classes
of maps form a quotient category
of $\text{Top}$, the {\it homotopy category
$\text{hTop}$}, where composition
of homotopy classes is induced by composition
of representing maps. If $f:X\to
Y$ represents an isomorphism in
$\text{hTop}$, then $f$ is called a
{\it homotopy equivalence} or
{\it $\text{h}$-equivalence}. In explicit
terms this means $f:X\to
Y$ is a homotopy equivalence if there exists
$g:Y\to X$, a
{\it homotopy inverse of $f$}, such that $gf$
and $fg$ are both
homotopic to the identity. Spaces $X$ and $Y$
are called
{\it homotopy equivalent} or of the same {\it
homotopy
type} if there exists a homotopy equivalence
$X\to Y$. A space is
{\it contractible} if it is homotopy
equivalent to a point. A map
$f:X\to Y$ is {\it null homotopic} if it is
homotopic to a constant
map.

\end{itemize}

\section{n-th homotopy
group of a pointed space}
\label{section-n-th-homot-group-point-spac}

\begin{definition}
\label{definition-homotopy-group}
Let $(X,x_0)$ be a pointed topological space
and $s_0\in S^n$. The elements of
the {\it $n$-th homotopy group} are homotopy
classes of maps
$(S^n,s_0)\to (X,x_0)$. Equivalently, they
are homotopy classes of maps
$(I^n,\partial I^n)\to (X,x_0)$. (Homotopies
are required to preserve the base
points, $s_0\mapsto x_0$ or $\partial
I^n\mapsto x_0$.)

Also,
$$
\pi_n(X,*)=[(I^n,\partial
I^n),(X,\{*\})]\cong[I^n/\partial I^n,X]^0
$$
where $[X,Y]$ denotes the set of homotopy
classes $[f]$ of maps $[f]:X\to Y$.
\end{definition}

\section{n-th relative
homotopy group of
a pointed pair of spaces}
\label{section-relative-homotopy-group}

\noindent
Let $A$ be a subspace of $X$ and $x_0\in A$.
The elements of the {\it relative homotopy
group $\pi_n(X,A,x_0)$}
are homotopy classes of maps $(I^n,\partial
I^n,J^{n-1})\to (X,A,x_0)$
where $J^{n-1}$ is the union of all but one
face of $I^n$.
That is,
$$
\pi_{n+1}(X,A,*)=[(I^{n+1},\partial
I^{n+1},J^n),(X,A,x_0)].
$$

The elements of such a group are homotopy
classes of based maps
$D^n\to X$ which carry the boundary $S^{n-1}$
into $A$.
Two maps $f,g$ are called {\it homotopic
relative to $A$}
if they are homotopic by a
basepoint-preserving homotopy
 $F:D_n\times[0,1]\to X$ such that, for each
 $p$ in $S^{n-1}$ and $t$ in $[0,1]$,
the element $F(p,t)$ is in $A$. Ordinary
homotopy groups are
recovered for the case in which $A=\{x_0\}$.

\begin{remark}
\label{remark-relative-homotopy}
This construction is motivated by studying the kernel of the
map induced by inclusion
\[
i_*:\pi_n(A,x_0)\longrightarrow\pi_n(X,x_0).
\]
\emph{[The rest of this lecture remark was not recorded.]}
\end{remark}

\section{Commutativity of higher homotopy
groups}
\label{section-commu-highe-homot-group}

\begin{proposition}
\label{proposition-commu-highe-homot-group}
$\pi_n(X,x_0)$ is an abelian group for $n>1$.
\end{proposition}

\section{Action of fundamental groupoid
\texorpdfstring{$\Pi_1(X)$}{Π₁(X)} of
\texorpdfstring{$X$}{X} on
\texorpdfstring{$\pi_{n}(X)$}{π₁(X)}}
\label{section-actio-funda-group}

\begin{itemize}
\item A {\it groupoid} is a small category in
which every morphism is an isomorphism.
\item The {\it fundamental groupoid}
$\Pi_{1}(X)$
is the category whose objects
are the points of  $X$ and
morphisms are homotopy classes of paths
joining two points.
Composition is defined doubling the speed.
\end{itemize}

The fundamental groupoid acts on homotopy
group functorially.

\section{Homotopy groups of a product of
spaces
are products of homotopy groups of factors}
\label{section-homot-group-produ-space}

\section{Higher homotopy groups of a covering
space
are the same as those of a base space}
\label{section-highe-homot-group-cover}

\section{Long exact sequence of homotopy
groups associated to a pair}
\label{section-long-exact-seque-homot-grou}

For any pair $(X,A,x)$ we have a long exact
sequence
$$
\xymatrix@C=1.5em{
\pi_n(A,x_0)\ar[r]^{i_*}
& \pi_{n}(X,x_0)\ar[r]^{j_*}
& \pi_{n-1}(X,A,x_0)\ar[r]^{\partial}
& \pi_{n-1}(A,x_0)\ar[r]
& \cdots\ar[r]
& \pi_0(X,x_0)
}
$$
where $i$ and $j$ are the inclusions
$(A,x_0)\hookrightarrow(X,x_0)$ and
$(X,x_0,x_0)\hookrightarrow(X,A,x_0)$. The
map $\partial$ comes from restricting maps
$(I^n,\partial I^n,J^{n-1})\to (X,A,x_0)$ to
$I^{n-1}$, or by restricting maps
$(D^n,S^{n-1},s_0)\to (X,A,x_0)$. The map,
called the {\it boundary map}, is a
homomorphism when $n>1$.

\section{n-connected
space}
\label{section-n-connected-space}

A space $X$ with basepoint $x_0$ is called
{\it $n$-connected} if $\pi_i(X,x_0)=0$ for
$i\leq n$. Thus 0-connected means
path-connected and 1 connected means
simply-connected.

\section{n-connected
pair}
\label{section-n-connected-pair}

A pair $(X,A)$ is {\it $n$-connected} if
$\pi(X,A,x_0)=0$ for $i\leq n$.

\section{n-equivalence of spaces}
\label{section-n-equivalence-of-spaces}

A map $f:X\to Y$ is an {\it  $n$-equivalence}
if the induced map on homotopy groups
$f_{*}:\pi_{i}(X)\to \pi_{i}(Y)$ is an
isomorphism for $i<n$ and a surjection for
$i=n$.

\begin{remark}
\label{remark-what-i-picke-up-lectu-if}
\begin{itemize}
\item $(X,A)$ is $n$-connected if and only if
$A\hookrightarrow X$ is an $n$-equivalence.
\item $f:X\to Y$ is an $n$-equivalence if and
only if when we replace $f$ by a cofibration
$X\to M_{f}$, $(M_{f},X)$ is $n$-connected.
\end{itemize}
\end{remark}

\section{Compactly generated spaces}
\label{section-compactly-generated-spaces}

(What I picked up from the lecture.) If $X$
is compactly generated, then $X$ is
weakly Hausdorff if the diagonal subset
$\Delta_X\subset X\times X$ is
$k$-closed.

Now some other notes from the lectures:

In $\text{CGWH}$, $\Hom(X,Y)$ is a space with
the compact-open topology.
This is a compactly generated space,
$\mathbf{k}(\Hom(X,Y))$.

(Also see
\href{https://en.wikipedia.org/wiki/%
Currying#Function_spaces}{wiki
on currying})
$$
\begin{aligned}
\text{Map}(X,Y):=\text{ the space of maps
}X\to Y.\\
\text{Map}(X\times
Y,Z)\cong\text{Map}(X,\text{Map}(Y,Z))\\
\Hom(X\times Y,Z)\cong
\Hom(X,\text{Map}(Y,Z))
\end{aligned}
$$
In the last line, product is product in
$\text{CGWH}$, not in $\text{Top}$.
The functor $-\times Y$ is left adjoint to
$\text{Map}(Y,-)$.

We don't care so much about $\text{Top}$. We
care much more about $\text{CGWH}$,
the full subcategory of $\mathit{Top}$ on
 {\it compactly generated weakly Hausdorff}
 spaces.

\begin{definition}[Is this right?]
\label{definition-is-this-right}
$X$ is {\it compactly generated} if, for any
subset $C\subset X$, and for all continuous
maps $f:K\to X$ from compact Hausdorff
spaces, $$\text{if } f^{-1}(C) \text{is
closed in }K\text{, then } C\text{ is
closed}.$$
\end{definition}
{\bf Claim.}
If $X$ is compactly generated, then $X$ is
weakly Hausdorff if the diagonal subset
$\Delta_X\subset X\times X$ is $k$-closed.


{\bf Explanation from the book of May}
The ordinary category of spaces
allows pathology that obstructs a clean
development of the foundations.
The homotopy and homology groups of spaces
are supported on compact subspaces, and
it turns out that if one assumes a separation
property
that is a little weaker
than the Hausdorff property,
then one can refine the point-set topology of
spaces to eliminate
such pathology without changing these
invariants.

One major source of point-set level pathology
can be passage to quotient spaces.
Use of compactly generated topologies
alleviates this.

\begin{proposition}
\label{proposition-if-x-compa-gener-x-y}
If $X$ is compactly generated and $\pi:X\to
Y$ is a quotient map,
then $Y$ is compactly generated if and only
if $(\pi\times \pi)^{-1}(\Delta Y)$
 is closed in $X\times X$.
\end{proposition}

\noindent
The interpretation is that a quotient space
of a
compactly generated space by a “closed
equivalence relation”
is compactly generated.

\subsubsection{Weakly Hausdorff space}
\label{subsection-weakly-housdorff-space}

\subsection{Cartesian closedness of the
category $CGWH$}
\label{subsection-carte-close-categ-cgwh}
See "Compactly generated spaces" by Charles
Rezk.

\section{Wedge}
\label{section-wedge}

\section{Smash product of spaces, its
associativity in
\texorpdfstring{$\text{CGWH}$}{CGWH}}
\label{section-smash-produ-space-its-assoc}

The appropriate analogue of the Cartesian
product in the category of based spaces is
the {\it smash product} $X\wedge Y$ defined
by
$$X\wedge Y=X\times Y/X\vee Y.$$
Here $X\vee Y$ is viewed as the subspace of
$X\times Y$ consisting of those pairs $(x,y)$
such that either $x$ is the basepoint of $X$
or $y$ is the basepoint of $Y$.

\begin{example}[Basic example]
\label{example-basic-example}
	The smash product of two circles is a
	sphere, as may be seen by collapsing the
	generating circles on a revolution torus
	(see
	\href{https://en.wikipedia.org/wiki/%
	Smash_product#/media/File:
	Visualization_of_the_smash_product_of_%
	two_circles.gif}{Wikipedia}).
\end{example}

\section{Reduced and unreduced suspension}
\label{section-reduc-unred-suspe}
We also have the {\it suspension of pointed
spaces}, which is like usual suspension but
also collapsing the distinguished point,
which has become an interval:
$$\Sigma X=(I\times X)/(t,x)\sim (0,y)\sim
(1,y)\;\forall y\in X.$$

\section{Path and loop spaces of a pointed
space}
\label{section-path-loop-space-point-space}
[This is defined in Hatcher close to the
homotopy fiber construction--- it's related
to Blakers-Massey, Freudenthal theorem.]

\begin{definition}[Path fibration]
\label{definition-path-fibration}
An important case of the homotopy fiber
construction
is when the given function $f$ is the
inclusion of the basepoint $b_0$ into $B$.
Then $E_{f}$ is the space $PB$ of paths in
$B$ starting at $b_{0}$,
and $p:PB\to B$ sends each path to its
endpoint.
The fiber $p^{-1}(b_{0})$ is the loopspace
$\Omega B$
consisting of all loops in  $B$ based at
$b_{0}$.
Since $PB$ is contractible by progressively
truncating paths,
the long exact sequence of homotopy groups
for the path
fibration yields (another proof that)
$$
\pi_{n}(X,x_{0})\cong \pi_{n-1}(\Omega
X,x_{0})\qquad \forall n.
$$
\end{definition}

\section{For pointed spaces: smash with any
\texorpdfstring{$Y$}{Y}
is adjoint to
\texorpdfstring{$\text{Map}(Y,-)$}{Map(Y,-)}}
\label{section-point-space-smash-any-y}

[From
\href{https://en.wikipedia.org/wiki/%
Smash_product#}{Wikipedia}]
Adjoint functors make the analogy between the
tensor product
and the smash product more
precise. In the category of $R$-modules over
a commutative ring $R$, the tensor
functor $(-\otimes_{A} A)$ is left adjoint to
the internal Hom functor
$\Hom(A,-)$, so that
$$
\Hom(X\otimes A,Y)\cong\Hom(X,\Hom(A,Y)).
$$

In the category of pointed spaces, the smash
product plays the role the tensor product in
this formula: if $A$, $X$ are compact
Hausdorff then we have an adjunction
$$\text{Maps}_{*}(X\wedge A,Y)\cong
\text{Maps}_{*}(X,\text{Maps}_{*}(A,Y))$$
where $\text{Maps}_{*}$ denotes continuous
maps that send basepoint to basepoint, and
$\text{Maps}_{*}(A,Y)$ carries the
compact-open topology.

\section{Adjointness of Sigma and Omega}
\label{section-adjointness-of-and}
In particular, taking $A$ to be the unit
circle $S^{1}$, we see that the reduced
suspension functor $\Sigma$ is left adjoint
to the loop space functor $\Omega$:
$$\text{Maps}_{*}(\Sigma X,Y)\cong
\text{Maps}_{*}(X,\Omega Y).$$

[From lectures:] Then we have
$$\Hom_{\text{CGWH}_*}(\Sigma X,\Sigma
X)\cong\Hom_{\text{CGWH}_*}(X,\Omega\Sigma
X)$$
where $\Sigma X=S\wedge X$ and $\Omega\Sigma
X=\text{Map}(S^1,\Sigma X)$. That is,
$S^1\wedge-$ is adjoint to
$\text{Map}(S^1,-)$.

\section{Hurewicz homomorphism}
\label{section-Hurewicz-homomorphism}

\begin{definition}[Hurewicz homomorphism]
\label{definition-hurewicz-homomorphism}
$$
\begin{aligned}
h: \pi_{n}(X) &\longrightarrow H_{n}(X) \\
[f] &\longmapsto f_{*}(\alpha)
.\end{aligned}
$$
\end{definition}
where $\alpha$ is a generator of $
H_{n}(S^{n})$.

\section{Absolute and relative versions of
Hurewicz theorem}
\label{section-absol-relat-hurew-theor}

\begin{theorem}[Hurewicz]
\label{theorem-Hurewicz-absolute}
If $X$ is $n$-connected for some $n \geq 2$,
them $\tilde{H}_i(X)=0$ for $i<n$
and $\pi_{n}(X)\cong H_n(X)$.
\end{theorem}

\begin{lemma}
\label{lemma-hurew-homom-gives-homot-equiv}
If $X$ and $Y$ are simply connected CW-spaces
and  $f:X \to Y$ induces
isomorphisms $H_n(X) \cong H_n(Y)$ for all
$n$, then it is an homotopy
equivalence.
\end{lemma}

\begin{theorem}[Hurewicz]
\label{theorem-Hurewicz-relative}
If $(X,A)$ is $(n-1)$-connected, i.e.
$\pi_{n}(X,A)=0$ for $i \leq n-1$, for
some $n\geq 2$ and $A$ is simply connected,
then $H_i(X,A)=0$ for $i\leq n-1$
and $\pi_{n}(X,A)\cong H_n(X,A)$.
\end{theorem}
Perhaps $X$ should be $(n-1)$-connected for
the relative version too---
they are stated together in Hatcher.

\section{Weak factorization systems}
\label{section-weak-factorization-systems}

Here we follow \cite{Riehl}.

\begin{definition}
\label{definition-weak-factorization-system}
A {\it weak factorization system
$(\mathcal{L},\mathcal{R})$} on a category
$\mathcal{M}$ is comprised of two clases of
morphisms $\mathcal{L}$
and $\mathcal{R}$ so that

\begin{enumerate}
\item Every morphism in $\mathcal{M}$ may be
factored as a morphism in $\mathcal{L}$
followed by a morphism in $\mathcal{R}$:
$$
\xymatrix{
\bullet\ar[rr]^{f}\ar[rd]_{\mathcal{L}
\ni\ell}
&& \bullet\\
& \bullet\ar[ru]_{r\in\mathcal{R}}
}
$$

\item The maps in $\mathcal{L}$ have the {\it
left lifting property} with respect to each
map in $\mathcal{R}$ and equivalently the
maps in $\mathcal{R}$ have the {\it right
lifting property} with respect to each map in
$\mathcal{L}$, that is, any commutative
square
$$
\xymatrix{
\bullet\ar[d]_{\mathcal{L}\ni\ell}\ar[r]
& \bullet\ar[d]^{r\in\mathcal{R}}\\
\bullet\ar@{-->}[ur]\ar[r]
& \bullet
}
$$
admits a diagonal filler as indicated making
both triangles commute.
When this lift is unique,
we say the factorization system is {\it
orthogonal}.

\item The classes $\mathcal{L}$ and
$\mathcal{R}$ are each closed under
retracts in the arrow category:
given a commutative diagram
$$
\xymatrix{
\bullet\ar[r]\ar[d]_{t}\ar@/^1pc/@{-}[rr]
\ar@/^2pc/@{-}[rr]
& \bullet\ar[r]\ar[d]^{s}
& \bullet\ar[d]^{t}\\
\bullet\ar[r]\ar@/_1pc/@{-}[rr]\ar@/_2pc/@{-}
[rr]
& \bullet\ar[r]
& \bullet
}
$$

\noindent
if $s$ is in that class then so is its
retract $t$.
\end{enumerate}
\end{definition}

\begin{exercise}[3.1.8, Riehl]
\label{exercise-3-1-8-rihel}
Verify that the class of morphisms
$\mathcal{L}$ characterized
by the left lifting property against a fixed
class of morphisms $\mathcal{R}$ is closed
under coproducts,
closed under retracts, and contains the
isomorphisms.
\end{exercise}

\section{Model category structure on a
category}
\label{section-model-category-structure}

\begin{definition}[Lecture]
\label{definition-lecture}
A {\it model structure} on a category
$\mathcal{A}$
is a choice of subcategories
$\mathcal{W},\mathcal{C},\mathcal{F}$
called {\it weak-equivalences}, {\it
cofibrations} and {\it fibrations}
with the following properties:

\begin{enumerate}
\item[0.] All (finite) small limits an
colimits.

\item \textbf{(2 of 3)} Given
$\bullet\overset{f}{\to}\bullet\overset{g}
{\to}\bullet$,
if either 2 out of 3 among $f,g,f\circ g$
are in $\mathcal{W}$ then all of them are.

\item
$(\mathcal{C}\cap\mathcal{W},\mathcal{F})$
and
$(\mathcal{C},\mathcal{F}\cap\mathcal{W})$
are both weak factorization systems.

$(\mathcal{B},\mathcal{D})$ is a weak
factorization system. That is,

\begin{enumerate}
\item Any morphism in $\mathcal{A}$ can be
factored
as a morphism in $\mathcal{B}$
followed by a morphism in $\mathcal{D}$.

\item Lifts:

$$
\xymatrix{
\bullet\ar[d]_{f}\ar[r]&\bullet\ar[d]^{g}\\
\bullet\ar@{-->}[ur]^{\exists}\ar[r]&\bullet
}
$$

\item[(c')] Notice that the axiom of retracts
is not necessary. $r'\in\mathcal{R}\iff$ it
satisfies (b) for all $\ell\in\mathcal{L}$.
\end{enumerate}
\end{enumerate}
\end{definition}

\begin{definition}
\label{definition-item}
\begin{itemize}
\item $X$ is {\it fibrant} if
$X\to\text{pt}$.
\item $X$ is {\it cofibrant} if
$$
\xymatrix{
X\ar@{->>}[r]&X
}
$$
\item $X$ is {\it bifibrant} if
$$
\xymatrix{
0\ar@{^{(}->}[r]&X\ar@{->>}[r]&\text{pt}
}
$$
\end{itemize}
\end{definition}

\medskip\noindent
Two examples of model structures
are Str\o m and Quillen model structures.

\begin{example}[Two interesting model
category structures on $\text{CGWH}$]
\label{example-two-inter-model-categ-struc}
\begin{enumerate}
\item {\it Hurewicz-Str\o m model structure}.
\begin{itemize}
\item Cofibrations = Huerwicz cofibrations.
\item Fibrations = Hurewicz fibrations.
\item Weak equivalences = homotopy
equivalences.
\end{itemize}
\item {\it Quillen model structure}. Defined
on $\text{Top}$.
\begin{itemize}
\item Cofibrations = "retracts of relative
cell complexes".
\item Fibrations = Serre fibrations.
\item Weak equivalences = weak homotopy
equivalences.
\end{itemize}
\end{enumerate}
\end{example}

\subsection{Hurewicz fibrations and
cofibrations}
\label{subsection-hurew-fibra-cofib}

\section{Fibrations}
\label{section-fibrations}

\begin{definition}[Hatcher]
\label{definition-hatcher}
A map $p:E\to B$ is said to have the {\it
homotopy lifting property} with respect to a
space $X$ if, given a homotopy $g_t:X\to B$
and a map $\tilde{g}_0:X\to E$ lifting $g_0$,
so $p\tilde{g}_0=g_0$, then there exists a
homotopy $\tilde{g}_t:X\to E$ lifting $g_t$.

The {\it lift extension property for a pair
$(Z,A)$} asserts that every map $X\to B$ has
a lift $Z\to E$ extending a given lift
defined on the subspace $A\subset Z$. The
case $(Z,A)=(X\times I,X\times\{0\})$ is the
homotopy lifting property.

A {\it fibration} is a map $p:E\to B$ having
the homotopy property with respect to all
spaces $X$.
\end{definition}

\begin{definition}[\href{https://ncatlab.org%
/nlab/show/lift}{nLab}]
\label{definition-https-ncatl-org-nlab}
	A morphism $i$ has the {\it left lifting
	property with respect to a morphism $p$} and
	$p$ has the {\it right lifting property with
	respect to $i$} if for each morphisms $f$
	and $g$, if the outer square in the
	following diagram commutes, there exists
	$\phi$ (I think not necessarily unique)
	completing the diagram:
	$$
	\xymatrix@R=2.5em{
		A\ar[r]^{f}\ar[d]_{i}&X\ar[d]^{p}\\
		B\ar[r]_{g}\ar@{-->}[ur]^{\phi}&Y
	}
	$$
\end{definition}

\begin{definition}[\href{https://ncatlab.org%
/nlab/show/homotopy+lifting+property}{nLab}]
Let $C$ be a category with products and with
interval object $I$. A morphism $E\to B$ has
the {\it homotopy lifting property} if it has
the right lifting property with respect to
all morphisms of the form $(\text{id},0):Y\to
Y\times I$.
\label{definition-homot-lifti-prope-nlab}

	This means that for all commuting squares
	$$
	\xymatrix{
		Y\ar[r]^{f}\ar[d]&E\ar[d]^{p}\\
		Y\times I\ar[r]_{F}\ar@{-->}[ur]^{\sigma}&B
	}
	$$
	there exists a morphism $\sigma:Y\times I\to
	E$ such that both triangles in the former
	diagram commute.

	A {\it fibration} (also called {\it Hurewicz
	fibration}) is a mapping $p:E\to B$
	satisfying the homotopy lifting property for
	all spaces $X$.
\end{definition}

\subsubsection{Cofibrations}
\label{subsection-cofibrations}

\begin{definition}[Hatcher]
\label{definition-hatcher-2}
	If in the definition of a fibration as a map
	satisfying the homotopy lifting property we
	reverse the direction of the arrows, we
	obtain the dual notion of a {\it
	cofibration}. This is a map $i:A\to B$
	(which is $E\to B$ backwards) satisfying the
	following property: (we used to have an
	homotopy from some space $X$ into the base,
	now we have:) given an homotopy $g_{t}:A\to
	X$ and (a lift)  $\tilde{g}:B\to X$ such
	that $\tilde{g}_{0}i=g_{0}$, there exists a
	homotopy $\tilde{g}_{0}:B\to X$ such that
	$\tilde{g}_{t}i=g_{t}$.

	In the special case that $i$ is the
	inclusion of a subspace, this is the
	homotopy extension property, and the next
	proposition says that this is indeed the
	general case. So a cofibration is an
	inclusion satisfying the homotopy extension
	property.
\end{definition}

\begin{proposition}[4H.1]
\label{proposition-4h-1}
	If $i:A\to B$ is a cofibration, then $i$ is
	injective, and in fact a homeomorphism onto
	its image.
\end{proposition}

\begin{remark}[From Hatcher p. 461]
\label{remark-from-hatcher-p-461}
The duality between subobjects and quotient
objects is clear for abelian groups,
where subobjects are kernels and quotient
objects cokernels. The strict
topological analog of a kernel is a fiber of
a fibration. Dually, the
topological analog of a cokernel is the {\it
cofiber} $B/A$ of a
cofibration $A\hookrightarrow B$. If we make
an arbitrary map  $f:A\to B$ into a
cofibration $A\hookrightarrow M_{f}$, the
cofiber is the mapping cone
$C_{f}=M_{f}/(A\times {0}$.
\end{remark}

\begin{definition}[\href{https://%
en.wikipedia.org/wiki/
Homotopy_extension_property}{wiki}]
\label{definition-https-en-wikip-org-wiki}
	Let $X$ be a topological space and let
	$A\subset X$. We say that the pair $(X,A)$
	has the {\it homotopy extension property}
	if, given a homotopy $f_\bullet:A\to Y^I$
	and a map $\tilde{f}_0:X\to Y$ such that
		$$\tilde{f}_0\circ\iota=f_0$$\
		(so $\tilde{f}$ is the lift of $f_0:A\to
		Y$) then there exists an {\it extension} of
		$f_\bullet$ to a homotopy
		$\tilde{f}_\bullet:X\to Y^I$ such that
		$\tilde{f}_\bullet\circ\iota=f_\bullet$.

		That is,
		$$
		\xymatrix{
			A\ar[r]^{f_\bullet}\ar[d]_{\iota}
			&Y^I\ar[d]^{\pi_0}\\
			X\ar@{-->}[ur]^{\tilde{f}_\bullet}
			\ar[r]_{\tilde{f}_0}&Y
		}
		$$
		So there's some
		\href{https://en.wikipedia.org/wiki/%
		Currying#Function_spaces}{currying}
		to make usual homotopies $f_\bullet:A\times
		I\to Y$ look like $f_\bullet:A\to Y^I$. As
		said in our lectures: ``a homotopy $X\times
		I\to Y$ is the same as a map $X\to
		\text{Map}(I,Y)$".

	A map $\iota:A\to X$ is a {\it cofibration}
	if it has the homotopy extension property
	with respect to all topological spaces  $Y$.
\end{definition}

\subsubsection{Mapping cylinder}
\label{subsection-mapping-cylinder}

\begin{definition}[\href{https://ncatlab.org%
/nlab/show/mapping+cylinder}{nLab}]
\label{definition-https-ncatl-org-nlab-2}
Given a continuous map $f:X\to Y$ of
topological spaces, one can define its {\it
mapping cylinder} as a pushout
$$
\xymatrix{
		X\ar[r]^{f}\ar[d]&Y\ar[d]\\
		X\times
		I\ar[r]_{(\sigma_0)_*(f)}&\text{Cyl}(f)
	}
$$
	in $\text{Top}$, where $I=[0,1]$ and
	$\sigma:X\to X\times I$ is given by
	$x\mapsto(x,0)$.

	Set theoretically, the mapping cyllinder is
	usually represented as que quotient space
	$$\left(X\times I\coprod Y\right)\Big/\sim$$
	where $\sim$ is the smallest equivalence
	relation identifying $(x,0)\sim f(x)$ for
	all $x\in X$.
\end{definition}

\begin{definition}[\href{https://%
en.wikipedia.org/wiki/Mapping_cylinder}{wiki}
]
\label{definition-https-en-wikip-org-wik-2}
The {\it mapping cylinder} of a function $f$
between topological spaces $X$ and $Y$ is the
quotient
$$M_f=(([0,1]\times X)\amalg Y)\big/\sim$$
where $\amalg$ denotes disjoint union, and
$\sim$ is the equivalence relation generated
by
$$(0,x)\sim f(x)\text{ for each }x\in X.$$
{\it That is, the mapping cylinder $M_f$ is
obtained by gluing one end of $X\times[0,1]$
to $Y$ via the map $f$.} Notice that the
``top" of the cylinder $\{1\}\times X$ is
homeomorphic to $X$, while the ``bottom" is
the space $f(X)\subset Y$.

Looks the mapping cylinder is just deforming
$X$ to $Y$ putting $X$ inside $Y$ via $f$…
\end{definition}

\begin{definition}[From homework]
\label{definition-from-homework}
	Let $f:X\to Y$ be a map. Let
	$M_f=X\times[0,1]\cup_fY$ be the {\it
	mapping cylinder of $f$}, i.e. the pushout
	of
	$X\overset{\cong}{\to}X\times\{0\}
	\hookrightarrow
	X\times[0,1]$ and of $f:X\times Y$.
\end{definition}

\subsection{Serre fibrations}
\label{subsection-serre-fibrations}

\begin{definition}[Hatcher, p. 376]
\label{definition-hatcher-p-376}
		The map $p:E\to B$ is said to have the {\it
		homotopy lifting property for a pair}
		$(X,A)$ if each homotopy $f_t:X\to B$ lifts
		to a homotopy $\tilde{g}_t:X\to E$ starting
		with a given lift $\tilde{g}_0$ and
		extending a given lift $\tilde{g}_t:A\to
		E$. In other words, the homotopy lifting
		property for $(X,A)$ is the lift extension
		property for $(X\times I,X\times\{0\}\cup
		A\times I)$.

	The important observation here is that
	homotopy lifting property for disks is
	equivalent to the homotopy lifting property
	for all CW pairs $(X,A)$. A map $p:E\to B$
	satisfying the homotopy lifting property for
	disks is sometimes called a {\it Serre
	fibration}.
\end{definition}

\subsection{Fiber bundles}
\label{subsection-fiber-bundles}

Here I include some more notes of the
discussion following Serre fibrations on
Hatcher.

	A {\it fiber bundle} structure on a space
	$E$, with fiber $F$, consists of a
	projection map $p:E\to B$ such that each
	point $B$ has a neighbourhood $U$ for which
	there is a homeomorphism $h:p^{-1}(U)\to
	U\times F$ making the following diagram
	commute
	$$
	\xymatrix{
		p^{-1}(U)\ar[rr]^{h}\ar[dr]_{p}
		&&U\times F\ar[dl]\\
		&U
	}
	$$
\begin{example}
\label{example-proje-space-yield-inter}
Projective spaces yield interesting fiber
bundles. In the real case we have the
familiar covering spaces $S^n\to\mathbb{R}
P^n$ with fiber $S^0$. Over the complex
numbers the analog of this is a fiber bundle
$S^1\to S^{2n+1}\to\mathbb{C} P^n$. Here
$S^{2n+1}$ is the unit sphere in
$\mathbb{C}^{n+1}$ and $\mathbb{C} P^n$ is
viewed as the quotient space of $S^{2n+1}$
under the equivalence relation
$(z_0,\ldots,z_n)\sim\lambda(z_0,\ldots,z_n)$
for $\lambda\in S^1$. The projection
$p:S^{2n+1}\to\mathbb{C} P^n$ sends
$(z_0,\ldots,z_n)$ to its equivalence class
$[z_0,\ldots,z_n]$.

To see that the local triviality condition
for fibre bundles is satisfied, …

	The constructino of the bundle $S^1\to
	S^{2n+1}\to\mathbb{C} P^n$ also works when
	$n=\infty$, so there is a fiber bundle
	$S^1\to S^\infty\to \mathbb{C} P^\infty$.

	The case $n=1$ is particularly interesting
	since $\mathbb{C} P^1=S^2$ and bundle
	becomes $S^1\to S^3\to S^2$ with fiber,
	total space, and base all speres. This is
	known as the {\it Hopf bundle}. The
	projection $S^3\to S^2$ can be taken to be
	$(z_0,z_1)\mapsto
	z_0/z_1\in\mathbb{C}\cup\{\infty\}=S^2$.
	(That is, seeing $S^2$ as the one-point
	compactification of $\mathbb{C}$.)

	In polar coordinates we may see $S^3$ as the
	union of several tori. Stereorgraphic
	projection yields the following figure:
	% Figure omitted: hopf-bundle.
	The limiting cases $T_0$ and $T_\infty$
	correspond to the unit circle in the
	$xy$-plane and the $z$-axis under the
	stereographic projection.	Each torus
	$T_\rho$ is aunion of circle fibers. These
	fiber circles have slope 1 on the torus,
	winding around once longitudinally and once
	meridionally. As $\rho$ goes to $0$ or
	$\infty$ the fiber circles approach the
	circles $T_0$ and $T_\infty$, which are also
	fibers. The figure below shows four tori
	decomposed into fibers:
	% Figure omitted: hopf-bundle2.

	Replacing the field $\mathbb{C}$ by the
	quaternions $\mathbb{H}$, the same
	constructions yield fiber bundles $S^3\to
	S^{4n+3}\to\mathbb{H} P^{n}$ over
	quaternionic projective spaces $\mathbb{H}
	P^n$. Here the fiber $S^3$ is the unit
	quaternions, and $S^{4n+3}$ is the unit
	sphere in $\mathbb{H}^{n+1}$. Taking $n=1$
	gives a second Hopf bundle $S^3\to S^7\to
	S^4=\mathbb{H} P^1$.

	Another Hopf bundle $S^7\to S^{15}\to S^8$…
\end{example}

\subsection{Homotopy equivalences and weak
homotopy equivalences}
\label{subsection-homot-equiv-weak-homot}

\begin{definition}[Weak homotopy equivalence]
\label{definition-weak-homotopy-equivalence}
	Any map $f:X\to Y$ that induces isomorphisms
	on all homotopy groups ( $\forall n\geq 0$
	and all choices of basepoint in $X$).
\end{definition}

\subsection{Any Hurewicz fibration is Serre
fibration}
\label{subsection-any-hurew-fibra-serre}
\subsection{Any relative CW-complex is a
cofibration in Quillen model structure}
\label{subsection-any-relat-cw-compl-cofib}
\subsection{Any map can be replaced by…}
\label{subsection-any-map-can-be-replaced-by}
\subsubsection{an homotopy equivalence
followed by Hurewicz fibration}
\label{subsection-homot-equiv-follo-hurew}
\subsubsection{a Hurewicz cofibration
followed by homotopy equivalence}
\label{subsection-hurew-cofib-follo-homot}
\subsubsection{a Quillen cofibration followed
by weak equivalence}
\label{subsection-quill-cofib-follo-weak}
This one follows from CW-approximation
theorem

\section{Homotopy theory of spaces (B)}
\label{section-homotopy-theory-of-spaces-b}

\subsection{Blakers-Massey theorem}
\label{subsection-blakers-massey-theorem}

As in Hatcher and in terms of homotopy
pullbacks and pushouts. (Be careful if you
refer to REsk's paper for the latter version,
since in his paper the $n$ in definition of
$n$-connectiviy is shifted by $1$.
\begin{theorem}[Blakers-Massey. Hatcher 4.23,
Tom Dieck 6.4.1]
\label{theorem-blake-masse-hatch-4-23-tom}
	Let $X$ be a CW complex decomposed as the
	union of subcomplexes $A$ and $B$ with
	nonempty connected intersection $C=A\cap B$
	(that is also a CW complex?). If $(A,C)$ is
	$m$-connected and $(B,C)$ is $n$-connected,
	$m,n\geq 0$, then the map
	$$\pi_{i}(A,C)\to \pi_{i}(X,B)$$
is an isomorphism for $i<m+n$ and a
surjection for $i=m+n$.
\end{theorem}

\begin{proposition}[Quotient theorem.
Hatcher, 4.28and Tom Dieck 6.10.2]
\label{proposition-quoti-theor-hatch-4}
	If a CW pair $(X,A)$ is $r$-connected and
	$A$ is $r,s \geq 0$, then the map
	$\pi_{i}(X,A)\to \pi_{n}(X/A)$ induced by
	the quotient map $X\to X/A$ is an
	isomorphism for $i\leq r+s$ and a surjection
	for $i=r+s+1$.
\end{proposition}
\begin{proof}
	Using cones.
\end{proof}

\subsection{Freudenthal suspension theorem}
\label{subsection-freud-suspe-theor}

\begin{lemma}[Freudenthal suspension theorem,
Hatcher 4.24]
\label{lemma-freud-suspe-theor-hatch-4}
	The suspension map $\pi_{i}(S^{n})\to
	\pi_{i+1}(S^{n+1})$ is an isomorphism for
	$i<2n-1$ and a surjection for $i=2n-1$. More
	generally, this holds for the suspension
	$\pi_{i}(X)\to \pi_{i+1}(SX)$ whenever $X$
	is an $(n-1)$-connected CW complex.
\end{lemma}
\begin{theorem}[Freudenthal suspension
theorem, lecture notes]
\label{theorem-freud-suspe-theor-lectu}
		Let $X$ be $\ell$-connected space. Then
		$X\to\Omega\Sigma X$ is a
		$(2\ell+1)$-equivalence, that is,
	$$\pi_i(X)\to \pi_{i+1}(\Sigma X),$$
	is a bijection for $i<2\ell+1$ and a
	surjection for $i=2\ell+1$..
\end{theorem}
\begin{proof}[Proof 2 from lectures]
	Consider
	$$
	\xymatrix{
		X\ar[r]\ar[d]&CX\ar[d]\\
		CX\ar[r]&\Sigma X
	}
	$$
	Then we use relative homotopy long exact
	sequence with $(X,CX)$ to get
	$\pi_i(CX,X)\cong\pi_{i-i}(X)$, which is
	zero for $0\leq i\leq \ell+1$. Then use
	relative homotopy exact sequence for the
	pair $(\Sigma X,CX)$. then we get that
	$\pi_i(\Sigma X,CX)=\pi_i(\Sigma X)$. And
	then if you use it for $(\Sigma X, X)$ and

	But it also turns out that $\pi_i(\Sigma
	X)=\pi_{i-1}(\Omega\Sigma X)$ because
	$$
\begin{aligned}
		\pi_n(Z)&=\Hom_{\text{hTop},*}(S^n,Z)=
		\Hom(S^1\wedge
		S^{n-1},Z)=\text{Hom}(S^{n-1},\Omega
		Z)=\pi_{n-1}(\Omega,Z)
	\end{aligned}
$$.
	And then since $CX\hookrightarrow \Sigma X$
	we get an arrow $\pi_i(CX,X)\to\pi_i(\Sigma
	X,CX)$ which is isomorphism for $0\leq i
	\leq 2\ell +1$ and surjective for
	$i=2\ell+2$.

	So apply Blakers-Massey an ell equalities to
	get maps fro
	$\pi_{i-1}(X)\to\pi_{i-1}(\Omega\Sigma X)$
	for $i$ as desired.
\end{proof}
\begin{lemma}
\label{lemma-if-x-conne-then-x-1-conne}
	If $X$ is $\ell$-connected, then $\Sigma X$
	is $(\ell+1)$-connected for $\ell\geq0$.
\end{lemma}
$$\begin{matrix}
	\text{space}&&S^0&\Sigma
	S^0=S^1&\Sigma^2S^0=S^2&\Sigma^3S^0=S^3&
	\cdots&\Sigma^nS^0=S^n\\
	\text{conectedness}&&-1&0&1&2&\cdots&(n-1)
\end{matrix}$$

\subsubsection{\texorpdfstring{$\pi_{i}(S^{n}
)=0$}{πᵢ(Sⁿ)}
for \texorpdfstring{$i<n$}{i<n}}
\label{subsection-s-for-i-n}

\begin{lemma}
\label{lemma-s-n-is-n-1-connected}
	$S^n$ is $(n-1)$-connected.
\end{lemma}
Back to Hopf fibration:
$$S^1\hookrightarrow S^3\to S^2$$
we get
$$0=\pi_2(S^3)\to\pi_2(S^2)\overset{\cong}
{\to}\pi_1(S^1)\to\pi_1(S^3)=0,$$
so
$$\mathbb{Z}=\pi_2(S^2).$$
Now consider a map $S^n\to S^n$. We get a map
$CS^n\to CS^n$ (in general, for $f:X\to Y$ we
have $(t,x)\mapsto(t,f(x))$ in the cones). We
also have $CS^n\to CS^n/S^n=S^{n+1}$.

Now if we take $\text{id}:S^n\to S^n$ we
shall get $\text{id}:S^{n+1}\to S^{n+1}$.
Think about this like
$\pi_1(S^1)\to\pi_2(S^2)$. Now from
Freudenthal we get
$\pi_{i-1}(X)\to\pi_i(\Sigma X)$ is
surjective because $i=0$. From Hopf fibration
we have $\pi_2(S^2)=\mathbb{Z}$. So we have a
surjective map $\mathbb{Z}\to\mathbb{Z}$. So
it's an isomorphism and we conclude that
$\text{id}_{S^2}$ is a generator of
$\pi_2(S^2)$.
\begin{lemma}
\label{lemma-since-s-n-n-1-conne-we-have}
	Since $S^n$ is $(n-1)$-connected, we have
	$$\pi_i(S^n)\to\pi_{i+1}(S^{n+1})$$
	is isomorphism for $i\leq 2(n-1)=2n-1$ and
	epimorphism form $i=2n-1$. (We just shift
	the indices of the Freudenthal theorem by
	one.)
\end{lemma}
\begin{lemma}
\label{lemma-n-s-n-z-id-s-n-as-gener}
	$\pi_n(S^n)=\mathbb{Z}$ with
	$\text{id}_{S^n}$ as generator.
\end{lemma}

\subsubsection{$\pi_{n}(S^{n})=\mathbb{Z}$}
\label{subsection-n-s-n-z}
\begin{lemma}
\label{lemma-n-k-s-n-n-k-1-s-n-1-isomo}
	$\pi_{n+k}(S^n)\to\pi_{n+k+1}(S^{n+1})$ is
	isomorphism for $k\leq n-1$ and epimorphism
	for $k=n-1$.
\end{lemma}
So for example
$$\pi_4(S^3)=\pi_5(S^4)=\pi_6(S^5).$$
And in fact they are $\mathbb{Z}/2$. This is
what are called the {\it $k$th stable
homotopy groups of a sphere}. And more in
general, we take any space and apply
$\Omega\Sigma$ enough times, and the homotopy
will start to stabilize.

Or for example from
$$S^1\hookrightarrow S^3\to S^2$$
we get
$$0=\pi_i(S^1)\to
\pi_i(S^3)\overset{\cong}{\to}\pi_i(S^2)
\to\pi_{i-1}(S^2)=0$$
So $\pi_3(S^2)\cong\mathbb{Z}$ in case you
were wondering.
{\bf Claim.}[Serre]
	$\pi_{4n-1}(S^{2n})\cong\mathbb{Z}
	\oplus$finite
	abelian. And $\pi_i(S^k)$ is finite abelian
	in all other cases.

\subsection{Whitehead theorem that weak
equivalence between CW-complexes is a
homotopy equivalence}
\label{subsection-white-theor-weak-equiv}

\begin{theorem}[Whitehead]
\label{theorem-whitehead}
	Weak homotopy equivalence for CW complexes
	is homotopy equivalence.

	An inclusion inducing isomorphism on all
	homotopy groups is a deformation retract.
\end{theorem}


and related theorem that weak equivalence
between bifibrant objects is a homotopy
equivalence.
\subsection{Cellular approximation theorem}
\label{subsection-celul-appro-theor}

\begin{definition}[Cellular map]
\label{definition-cellular-map}
	A map between CW complexes $f:X\to Y$ is
	called {\it cellular} if $f(X^{n})\subset
	Y^{n}$ for all $n$.
\end{definition}
\begin{theorem}
\label{theorem-every-map-betwe-cw-compl}
	Every map between CW complex is homotopic to
	a cellular map.
\end{theorem}
\begin{theorem}[CW approximation]
\label{theorem-cw-approximation}
	Every space has a CW approximation, ie. it
	is weakly equivalent to a CW-complex.
	If $X$ is path-connected, the CW-complex can
	be chosen to have a single 0-cell.

	If $(X,A)$ is an $n$-connected CW-pair (so
	$A$ is a sub-CW-complex of $X$, made up of
	cells from $X$) there exists a CW-pair
	$(Z,A)\simeq (X,A)\text{rel}A$ with cells in
	$Z\setminus A$ of dimension greater than
	$n$.
\end{theorem}


\section{Postnikov towers}
\label{section-postnikov-towers}
Let $X$ be a CW complex. There exist spaces
$$
\xymatrix{
&\vdots\ar[d]\\
&Y_2\ar[d]\\
&Y_1\ar[d]\\
X\ar[uur]\ar[ur]\ar[r]&Y_0
}
$$
that are like the original space
(I think this means that the arrows in the
above diagram induce isomorphisms in
homotopy groups for smaller indices)
but with trivial homotopy groups for $k\geq
n$.
They are obtained by attaching to the
previous one
(or to $X$ itself?)
cells of dimension $n+2$
(because this kills the $n+1$-th homotopy
group
so that only the $n$-th homotopy group
remains nontrivial).

The inclusions
$X\hookrightarrow Y_{n}\hookrightarrow
Y_{n-1}\hookrightarrow
\ldots\hookrightarrow Y_{1}=K(\pi_{1}(X),1)$
 are converted to fibrations,
and in fact the fiber of $Y_{k}\to Y_{k-1}$
is a $K(\pi_{k}(X),k)$-space.

\begin{remark}
\label{remark-postn-tower-some-kind-facto}
It looks like Postnikov towers posses an
interesting property
of being a factorization system for some
category.
See
\url{https://ncatlab.org/nlab/show/Postnikov%
+system}.
Recall that this means that the maps in that
category
will be factored through maps in the
subcategories of right
and left morphisms provided by the Postnikov
tower.
\end{remark}


\section{Whitehead towers}
\label{section-whitehead-towers}

Now following Laurentium Maxim,
"Lecture notes on homotopy theory and
applications".

\begin{definition}[\href{https://%
people.math.wisc.edu/~lmaxim/754notes.pdf}
{754notes.pdf}]
\label{definition-https-peopl-math-wisc}
A {\it Whitehead tower} of $X$ is a sequence
of fibrations
$$
\xymatrix{
\cdots\ar[r]&X_n\ar[r]&X_{n-1}\ar[r]&\cdots
\ar[r]&X_0=X
}
$$
such that
\begin{itemize}
\item $X_n$ is $n$-connected.
\item $\pi_{q}(X_{n})=\pi_{q}(X)$ for $q\geq
n+1$.
\item The fiber of $X_n\to  X_{n-1}$ is a
$K(\pi_{n}(X),n-1)$-space.
\end{itemize}
\end{definition}

\begin{lemma}
\label{lemma-white-tower-exist-cw-compl}
Whitehead towers exist for CW complexes.
\end{lemma}
\begin{proof}
	Consider
	$$\text{hofib}\hookrightarrow X\to
	\text{"$X$ without $\pi_{k}$, $k>n$}"$$
	We immediately get that the second arrow
	induces an isomorphism on $\pi_{i}$ for
	$i\leq n$, while the first arrow induces
	isomorphisms for $i>n$. Define $
	X_{n}:=\text{hofib}$.

	Now notice that $X_n$ and $X_{n+1}$ have
	isomorphic homotopy groups except perhaps
	the $(n+1)$-th, because for $X_{n}$ it is
	already the same as that of $X$, but for
	$X_{n+1}$ it is still zero. So, from the
	fibration
	$$\text{hofib} \hookrightarrow X_{n+1}\to
	X_{n}$$
	we get an isomorphism
	$$
	\xymatrix@C=.9em@R=.9em{
		\cdots\ar[r]
		&\pi_{n+2}(X_{n})\ar[r]\ar[d]_{\cong}
		&\pi_{n+1}(\text{hofib} )\ar[d]^{=}\ar[r]
		&\pi_{n+1}(X_{n+1})\ar[r]\ar[d]^{=}
		&\pi_{n+1}(X_{n})\ar[r]\ar[d]_{\cong}
		&\pi_{n}(\text{hofib} )\ar[r]
		&\pi_{n}(X_{n+1})\ar[r]\ar[d]^{=}
		&\cdots\\
		&\pi_{n+2}(X)&0&0
		&\pi_{n+1}(X)\ar[ur]_{\cong}
		&&0&
	}
	$$
	showing that the homotopy fiber of this
	fibration is in fact $K(\pi_{n+1}(X),n)$.

\end{proof}

\subsection{How to kill higher homotopy
groups without changing lower homotopy groups
and lower homology}
\label{subsection-how-kill-highe-homot}

Looks like Blakers-Massey is needed for this.
The following are lecture notes corresponding
to post-Blakers-Massey and pre-Postnikow
towers part our the course.

Glue a disk to a space and what happens to
homotopy groups?
$$
\xymatrix{
	S^{n-1}\ar[r]^{(n-1)\text{-equiv}}\ar[d]
	_{0\text{-equiv}}
	&D^n\ar[d]\\
	X\ar[r]&X\cup D^n\ar@{}[ul]|{\ulcorner}
}
$$
Assume $X$ is connected. We get a map from
the vertical arrows
$$
\xymatrix{
	\pi_i(D^n,S^{n-1})\ar[r]&\pi_i(X\cup D^n,X)
}
$$
which is $(n-1)$-equivalence by
Blakers-Massey. So, by attaching $\sqcup D^n$
we can kill $\pi_{n-1}(X)$, that is,
$X\cup(\sqcup D^n)$ has trivial $\pi_{n-1}$.

Now notice that

% since $S^{n-1}\to D^n$ is an
% $(n-1)$-equivalence, which is the case
% since
$$
\xymatrix{
	0=\pi_i(D^n)\ar[r]
	&\pi_i(D^n,S^{n-1})\ar[r]^{\cong}
	&\pi_{i-1}(S^{n-1})\ar[r]
	&\pi_{i-1}(D^n)=0
}
$$
that is, $\pi_i(D^n,S^{n-1})=0$ for $i\leq
n-1$. This implies that $\pi_i(X\cup
D^n,X)=0$ for $i\leq n-1$.

Also by homotopy long exact sequence,
$$\pi_{n-1}(X)\to \pi_i(X\cup D^n)\text{ is
sujrective}$$
$$\pi_i(X)\to \pi_{i}(X\cup D^n)\text{ is
isomorphism for }i\leq n-2.$$
So what we have thus far is
$$
\xymatrix{
	\pi_n(X\cup D^n)\ar[r]
	&\pi_{n-1}(X)\ar[r]
	&\pi_{n-1}(X\cup D^n)\ar[r]
	&0=\pi_{n-1}(X\cup D^n)
}
$$
Notice that $\pi_n(X\cup D^n,X)$ is not
ingeneral cyclic (counterexample $S^1\cup
D^2$ taking unieversal cover which is real
line with spheres on integers, homotopy
equivalent to $\bigvee_\mathbb{Z} S^2$ which
is not finitely generated).

\begin{quote}
	So basically attaching a disk via $f$ will
	kill $[f]$ inside $\pi_n(X)$ this is called
	{\it killing} an element of the homotopy
	group.
\end{quote}

\begin{proposition}
\label{proposition-any-cw-compl-x-x-x-equi}
	For any CW-complex $X$, $X^\ell\to X$ is an
	$\ell$-equivalence.
\end{proposition}

\begin{remark}
\label{remark-we-have-used-x-long-exact}
	We have used that for $A\hookrightarrow X$
	from long exact sequence of relative
	homotopy groups we get $\pi_n(X,A)=0$.
\end{remark}

\begin{lemma}
\label{lemma-let-i-g-d-i-x-g-d-i-x-then}
	Let $i\leq \ell$ and $g: D^i\to X$,
	$g:\partial D^i\to X^\ell$. Then there is a
	homotopy $\text{rel }\partial D^i$ to a map
	with $\text{img}\subset X^\ell$.
\end{lemma}

\begin{theorem}[Cellular approximation
theorem]
\label{theorem-cellu-appro-theor}
	Let $X$ and $Y$ be CW-complexes and
	$\xi:Y\to X$ be any map. Then $\xi$ is
	homotopic to a {\it cellular map}, that is,
	a map $\psi:Y\to X$, such that for all
	$\ell$, $\psi Y^0\subset X^\ell$.
\end{theorem}

We also have
\begin{proposition}
\label{proposition-let-n-2-then}
Let $n\geq 2$. Then
$$
\pi_n\left(
\bigvee_{k\in I}S^n\right)
=\bigoplus_{k\in I}\pi_n(S^n)
=\bigoplus_{k\in I}\mathbb{Z}=Z^{\oplus I}
$$
\end{proposition}

\begin{proposition}
\label{proposition-finite-i}
First notice that for finite $I$,
$$X^n=X^{n+1}=\vee_{k\in I}S^n$$
by looking carefully. Then
$$\pi_n(X,X^{n+1})=0=\pi_{n+1}(X,X^{n+1})$$
so
$$
\bigoplus_{k\in I}\mathbb{Z}
=\prod_{k\in I}\pi_n(S^n)
=\pi_n(X)=\pi_n(X^{n+1})
=\pi_n(X^n)=\pi_n\left(
\vee_{k\in I}S^n\right)
$$
and for the infinite case it also works,
using finite compactness of the CW complex.
\end{proposition}

\section{Homotopy theory of spaces (C)}
\label{section-homotopy-theory-of-spaces-c}

\subsection{Homotopy pushouts and pullbacks}
\label{subsection-homot-pusho-pullb}
The idea is simply to construct the homotopy
version of pushouts and pullbacks. So let's
remember that
\begin{definition}
\label{definition-item-2}
	\begin{itemize}
		\item A {\it pullback} of the morphisms $f$
		and $g$ consists of an object $P$ and two
		morphisms $p_1:P\to X$ and $p_2:P\to Y$
		satisfying the following universal
		property:
		$$
		\xymatrix{
			Q\ar@/^1pc/[rrd]^{q_2}\ar@/_1pc/[ddr]
			_{q_1}
			\ar@{-->}[dr]^{\phi}\\
			&P\ar[r]^{p_2}\ar@{}[rd]|{\lrcorner}
			\ar[d]_{p_1}&Y\ar[d]^{g}\\
			&X\ar[r]_{f}&Z
		}
		$$
		\item A {\it pushout} of the morphisms $f$
		and $g$ consists of an object $P$ and two
		morphisms $i_1:P\to X$ and $i_2:P\to Y$
		satisfying the following universal
		property:
		$$
		\xymatrix{
			Z\ar[r]^{g}\ar[d]_{f}
			&Y\ar[d]^{i_2}\ar@/^1pc/[ddr]^{j_2}\\
			X\ar[r]_{i_1}\ar@/_1pc/[drr]_{j_1}
			&P\ar@{}[ul]|{\ulcorner}\ar@{-->}[dr]
			^{\phi}\\
			&&Q
		}
		$$
	\end{itemize}
\end{definition}
\begin{remark}
\label{remark-other-names-pusho-cofib-prod}
Other names for the pushout are {\it
cofibered product of $X$ and $Y$} (especially
in algebraic categories when $i_1$ and $i_2$
are monomorphisms), or {\it free product of
$X$ and $Y$} with $Z$ {\it amalgamated sum},
or more simply an {\it amalgamation} or {\it
amalgam of $X$ and $Y$.}
\end{remark}
So basically
\begin{definition}[Homotopy pullback and
homotopy pushout]
\label{definition-homot-pullb-homot-pusho}
	We just ask that the diagrams above commute
	up to homotopy?
\end{definition}
\begin{remark}[Concrete construction of
homotopy pushout]
\label{remark-concr-const-homot-pusho}
	We also need to ask in our definition that
	there is some sort of homotopy-naturality.
	For the case of the homotopy pushout it
	should look something like this:
		$$
		\xymatrix{
			\bullet\ar[rr]\ar[dd]\ar[dr]
			&&\bullet\ar[d]^{\simeq}\\
			&A\ar[r]\ar[d]&B\ar[d]\ar@/^1pc/[ddr]\\
			\bullet\ar[r]^{\simeq}
			&C\ar[r]\ar@/_1pc/[drr]
			&\bullet\ar@{}[ul]|{\ulcorner}\ar[dr]
			^{\simeq}\\
			&&&\bullet
		}
		$$
	so everything is commuting up to homotopy.
	What are the formal definitions? Probably
	rather involved. But a concrete definition
	of the {\it homotopy pushout} according to
	the diagram above is
	$$B\sqcup C\sqcup A\times
	I\Big/\substack{A\times 0\sim B \\ A\times
	I\sim C}$$
\end{remark}
\begin{remark}[From lecture notes]
\label{remark-from-lecture-notes}
		The pullback of
	$$
	\xymatrix{
		&Z^I\times_ZY\ar[d]\\
		X\ar[r]&Z
	}
	$$
	is the homotopy pullback of
	$$
	\xymatrix{
		&Y\ar[d]\\
		X\ar[r]&Z
	}
	$$
\end{remark}
\begin{definition}[From lecture notes. Is
this right?]
\label{definition-lectu-notes-right}
	The {\it homotopy pullback} of a diagram
		$$
		\xymatrix{
			&Y\ar[d]\\
			X\ar[r]&Z
		}
		$$
		is
		$$
		\xymatrix{
			X\times_{\text{ev}_0}Z^I\times_{\text{ev}
			_1}Y
			\ar[d]\ar[r]&Y\ar[d]\\
			X\ar[r]&Z
		}
		$$
		Intuitively, for any $x\in X$ and $y\in Y$
		this object has the space of paths
		connecting $x$ and $y$.
\end{definition}

\subsection{Homotopy fibers and cofibers}
\label{subsection-homot-fiber-cofib}

Before we go, what is a fiber? Perhaps it is
the same as a kernel. And a cofiber may be
the same as a cokernel. Recall:
\begin{definition}[Kernel and cokernel]
\label{definition-kernel-and-cokernel}
	\begin{itemize}
		\item The {\it kernel} of a morphism is
		that part of its domain which is sent to
		zero. Formally, in a category with an
		initial object 0 and pullbacks, the {\it
		kernel $\ker f$} of a morphism $f:A\to B$
		is the pullback $\ker(f)\to A$ along $f$ of
		the unique morphism $0\to B$

		More explicitly, this characterizes the
		object $\ker(f)$ as \textit{the} object
		(unique up to isomorphism) that satisfies
		the following universal property:
		\begin{quote}
			for every object $C$ and every morphism
			$h:C\to A$ such that $f\circ h=0$ is the
			zero morphism, there is a unique morphism
			$\phi:C\to\ker(f)$ such that
			$h=p\circ\phi$.
		\end{quote}
		$$
		\xymatrix{
			C\ar@{-->}[dr]^{\phi}\ar@/_1pc/[ddr]_{h}\\
			&\ker(f)\ar[d]\ar[r]&0\ar[d]\\
			&A\ar[r]_{f}&B\ar@{}[ul]|{\ulcorner}
		}
		$$
		\item In a category with a terminal object
		1, the {\it cokernel} of a morphism $f:A\to
		B$ is the pushout (arrows $h$ and $\phi$
		apply if terminal object is zero)
		$$
		\text{coker}(f):=1\sqcup_AB
		\qquad\qquad
		\xymatrix{
			A\ar[r]^{f}\ar[d]\ar@{}[dr]|{\lrcorner}
			&B\ar[d]\ar@/^1pc/[rdd]^{h}\\
			1\ar[r]&\text{coker}(f)\ar@{-->}[dr]
			^{\phi}\\
			&&C
		}
		$$
		In the case when the terminal object is in
		fact zero object, one can, more explicitly,
		characterize the object $\text{coker}(f)$
		with the following universal property:
		\begin{quote}
			for every object $C$ and every morphism
			$h:B\to C$ such that $h\circ f=0$ is the
			zero morphism, there is a unique morphism
			$\phi:\text{coker}(f)\to C$ such that
			$h=\phi\circ i$.
		\end{quote}
\end{itemize}
\end{definition}
The {\it homotopy fiber} and the  {\it
homotopy cofiber} are probably just the
homotopic versions of the kernel and
cokernel.

\begin{definition}[From lecture notes, looks
incomplete…]
\label{definition-lectu-notes-looks-incom}
	The {\it homotopy fiber} if $f:Y\to Z$ is
	the pullback of
		$$
		\xymatrix{
			&Y\ar[d]^{f}\\
			pt\ar[r]&Z
		}
		$$
		$F\subset Z^I\times_ZY\to Z$, where $F$ is
		the space of paths starting at $x$ and
		ending at the same point $f(y)$.
\end{definition}

\begin{definition}[Homotopy fiber, a concrete
construction in Hatcher]
\label{definition-homot-fiber-concr-const}
	Given a map $f:A\to B$, let $E_f$ be the
	space of pairs $(a,\gamma)$ where $a\in A$
	and $\gamma:I\to B$ is a path in $B$ with
	$\gamma(0)=f(a)$. So: take a point in $A$,
	and map it to some point in $B$, say $f(a)$.
	And then take any path in $B$ starting at $f
	(a)$.
\begin{proposition}
\label{proposition-map-p-e-f-b-given-p}
	The map $p:E_f\to B$ given by
	$p(a,\gamma)=\gamma(1)$ is a fibration.
\end{proposition}
We can regard $A$ as a subspace of $E_{f}$
consisting of pairs $(a,\gamma)$ with
$\gamma$ the constant path at $f(a)$, and
$E_f$ deformation retracts onto this subspace
by restricting all the paths $\gamma$ to
shorter and shorter initial segments. The map
$p:E_f\to B$ restricts to $f$ on the subspace
$A$, so we have factored the map $f:A\to B$
as the composition $A\hookrightarrow E_f\to
B$ of a homotopy equivalence and a fibration.

We can also think of this construction as
extending $f$ to a fibration $E_f\to B$ by
enlarging its domain to a homotopy equivalent
space.

The fiber $F_f$ of $E_f\to B$ is called the
{\it homotopy fiber} of $f$. It consists of
the pairs $(a,\gamma)$ with $a\in A$ and
$\gamma$ a path in $B$ from $f(a)$ to a
basepoint $b_0\in B$.
\end{definition}

\subsection{Homotopy fiber of homotopy fiber}
\label{subsection-homot-fiber-homot-fiber}

It is interesting to see what happens when
the process of forming homotopy fibers is
iterated. Given a fibration $p:E\to B$ with
fiber $F=p^{-1} (b_0)$, we know that the
inclusion of $F$ into the homotopy fiber
$F_{p}$ is a homotopy equivalence (because
all fibers in homotopy fibrations are
homotopy equivalent?) Well, beacause
\begin{proposition}[4.66]
\label{proposition-4-66}
	If $p:E\to B$ is a fibration, then the
	inclusion $E\hookrightarrow E_p$ is a fiber
	homotopy equivalence (a fiber-preserving map
	of the total spaces that has an "fiber
	homotopy inverse" i.e. composition is
	homotopic to the identity via
	fiber-preserving maps). In particular, the
	homotopy fibers are homotopy equivalent to
	the actual fibers.
\end{proposition}

Recall that the homotopy fiber $F_p$ consists
of pairs $(e,\gamma)$ with $e\in E$ and
$\gamma$ a path in $B$ from $p(e)$ to $b_0$.
The inclusion $F\hookrightarrow E$ extends to
a map $i:F_{p}\to E$, $i(e,\gamma)=e$ and
\textit{this map is obviously a fibration}.
In fact it is the pullback via $p$ of the
path fibration $PB\to B$… what could this
mean? Maybe this:
$$
\xymatrix{
	F_{p}\ar[r]^{i}\ar[d]\ar@{}[dr]|{\lrcorner}
	&E\ar[d]^{p}\\
	PB\ar[r]&B
}
$$
Perhaps this justifies why, up to
equivalence, the fiber of this map $i$ must
be $\Omega B$. In the end we get the
following sequence:
$$
\xymatrix@C=1.5em{
	\cdots\ar[r]&\Omega^{2}B\ar[r]&\Omega
	F\ar[r]
	&\Omega E\ar[r]&\Omega
	B\ar[r]&F\ar[r]&E\ar[r]&B
}
$$
where any two consecutive maps form a
fibration, up to homotopy equivalence, and
all the maps to the left $\Omega B$ are
obtained by applying the functor $\Omega$ to
the later maps. The long exact sequence of
homotopy groups for any fibration in the
sequence coinincides with the long exact
sequence for $F\hookrightarrow E\to B$.

\subsection{Homotopy cofiber of homotopy
cofiber}
\label{subsection-homot-cofib-homot-cofib}

The dual construction to the sequence above
is
$$
\xymatrix@C=1.5em{
	X\ar[r]&Y\ar[r]&M_{f}\ar[r]&\Sigma X\ar[r]
	&\Sigma Y\ar[r]&\Sigma
	M_{f}\ar[r]&\Sigma^{2}X\ar[r]
	&\cdots
}
$$

\subsection{The two Dold-Puppe sequences,
their (co)-exactness}
\label{subsection-two-dold-puppe-seque}

\subsection{Fiber of a (Serre) Hurewicz
fibration is a (weak) homotopy fiber}
\label{subsection-fiber-serre-hurew-fibra}

\subsection{Fiber bundles are Serre
fibrations}
\label{subsection-fiber-bundl-serre-fibra}
We've been talking a lot about Hurewicz
fibrations. Let's talk about Serre
fibrations. Notice that H. fibration
$\implies$ S. fibration. What is the most
natural example of a Serre fibration?

\begin{proposition}[Hatcher, 4.48]
\label{proposition-hatcher-4-48}
	Let $E$ be a fiber bundle with fiber $F$.
	Then $f$ is a Serre fibration.
\end{proposition}
\begin{proof}
	What sdoes it mean to be a Serre fibration?
	It means that
	$$
	\xymatrix{
		I^n\ar[r]\ar[d]&E\ar[d]\\
		I^{n+1}=I^n\times I\ar[r]\ar@{-->}[ur]&B
	}
	$$
	So if $\mathcal{U}$ is a covering of $B$
	such that $f^{-1}U\cong U\times F$. By
	Lebesgue lemma, there is a $\delta>0$ such
	that for all $x\in I^{n+1}$, the ball
	$B(x,\delta)$ lies in some $f^{-1}U$ for
	some $U$.

	Then we subdivide $I^{n+1}$ in smaller cubes
	of the same size with diameter $<\delta$.
	So, each the image of each cube lies in some
	$U\in\mathcal{U}$.

	Then
	$$
	\xymatrix{
		I^n\ar[d]\ar[r]&F\times U\ar@{->>}[d]\\
		I^{n+1}\ar[r]\ar@{-->}[ur]&U
	}
	$$
	has a lift for every little square because
	$$
	\xymatrix{
		X\ar[r]\ar[d]&U\ar[d]\\
		X\times I\ar[r]\ar@{-->}[ur]&pt
	}
	$$
	is always a fibration (think about this) and
	because pullbacks of fibrations are
	fibrations:
	$$
	\xymatrix{
		U\times
		F\ar@{->>}[r]\ar@{->>}[d]&U\ar@{->>}[d]&\\
		F\ar@{->>}[r]&pt
	}
	$$.
	Then we may just add up the squares because
	$$
	\xymatrix{
		D^n\ar[d]\\
		D^n\times I
	}
	$$
	and we're done.
\end{proof}

\subsection{Long exact sequence of homotopy
groups associated to Serre fibration}
\label{subsection-long-exact-seque-homot}
\begin{theorem}[Hatcher 4.41]
\label{theorem-hatcher-4-41}
	Suppose $p:E\to B$ is a Serre fibration.
	Choose basepoints $b_{0}\in B$ and $x_{0}\in
	F=p^{-1}(b_{0})$. Then the map
	$p_{*}:\pi_{n}(E,F,x_{0})\to
	\pi_{n}(B,b_{0})$ is an isomorphism for all
	$n\geq 1$.

	Hence, if $B$ is path connected, there is a
	long exact sequence
$$
\xymatrix{
	\cdots\ar[r]&\pi_n(F)\ar[r]&\pi_n(E)\ar[r]
	&\pi_n(B)\ar[r]&\pi_{n-1}(F)\ar[r]
	&\pi_{n-1}(E)\ar[r]&\cdots\\
	\cdots\ar[r]&\pi_n(F)\ar[r]\ar[u]^{=}
	&\pi_n(E)\ar[r]\ar[u]^{=}
	&\pi_n(E,F)\ar[r]\ar[u]^{\cong}\ar[u]_{p_{*}
	}
	&\pi_{n-1}(F)\ar[r]\ar[u]^{=}
	&\pi_{n-1}(E)\ar[r]\ar[u]^{=}&\cdots
}
$$
\end{theorem}

\subsection{Fiber bundles over manifolds are
Hurewicz fibrations}
\label{subsection-fiber-bundl-over-manif}

\section{Homotopy theory of spaces (D)}
\label{section-homotopy-theory-of-spaces-d}

\subsection{Eilenberg-MacLane spaces}
\label{subsection-eilenberg-maclane-spaces}
A space $X$ having jut one nontrivial
homotopy group $\pi_{n}(X)\cong G$ is called
an {\it Eilenberg-MacLane space} $K(G,n)$. We
can build a CW complex $K(G,n)$ by the
familiar construction of attaching cells to a
wedge of spheres, acounting for generators
and relations, and then attaching further
cells to kill higher homotopy groups.

\begin{proposition}[4.30]
\label{proposition-4-30}
	The homotopy type of a CW complex $K(G,n)$
	is uniquely determined by $G$ and $n$.
\end{proposition}
\begin{proof}[Idea of proof]
	Attaching cells and Whithead theorem.
\end{proof}

\subsection{Representable functors and Yoneda
lemma}
\label{subsection-repre-funct-yoned-lemma}

\begin{definition}[Representable functor]
\label{definition-representable-functor}
	Let $\mathcal{C}$ be a category. A functor
	$F:\mathcal{C}^{\text{op}}\to \text{Sets}$
	is called {\it representable} if there
	exists an object $B=B_{F}$ in $\mathcal{C}$
	with the property that there is a
	\textit{natural} isomorphism of functors
	$$\varphi:\mathcal{C}(-,B_{F})\to F$$
	where $\mathcal{C}(-,B_{F})$ is the set of
	arrows from $-$ to $B_{F}$.

	One ussually expresses the naturality
	condition for a map $f:X\to Y$ with the
	following diagram:

	$$
	\xymatrix{
		\mathcal{C}(X,B)\ar[r]^{\varphi_{X}}\ar[d]
		_{f^{*}}
		&F(X)\ar[d]^{f^{*}}\\
		\mathcal{C}(Y,B)\ar[r]_{\varphi_{Y}}&F(Y)
	}
	$$
\end{definition}

\begin{remark}
\label{remark-apply-b-ident-map-b-b-we}
	By applying $\varphi_{B}$ to the identity
	map  $B\to B$ we obtain a special element
	$\gamma=\varphi_{B}(\text{id})\in F(B)$,
	which is \textit{generic} in the sense that
	any element $\xi \in F(X)$ can be obtained
	as $\xi=f^{*} (\gamma)$, for a suitable
	$f:X\to B$. Indeed, one can take
	$f=\varphi^{-1}_{X}(\xi)$ and apply
	naturality to check that
	$\xi=f^{*}(\gamma)$.
\end{remark}


\begin{lemma}[Yoneda,
\href{https://en.wikipedia.org/wiki/%
Yoneda_lemma}{wiki}]
\label{lemma-yoned-https-en-wikip-org-wiki}
	Let $F$ be a functor from a locally small
	category $\mathcal{C}$ to $\text{Set}$. Then
	for each object $X$ of $\mathcal{C}$, the
	natural transformations
	$\text{Nat}(\Hom(X,-),F)$ are in one-to-one
	correspondence with the elements of $F(X)$,
	that is
	$$\text{Nat}(\Hom(X,-), F)\cong F(X)$$
Moreover, this isomorphism is natural in $A$
and $F$ when both sides are regarded as
functors from
$\mathcal{C}\times\text{Set}^{\mathcal{C}}$
to $\text{Set}$. ($\text{Set}^{\mathcal{C}}$
denotes de category of functors from
$\mathcal{C}$ to $\text{Set}$.)

	There is a contravariant version of Yoneda
	lemma asserting that if $F$ is a
	contravariant functor from $\mathcal{C}$ to
	$\text{Set}$,
	$$\text{Nat}(\Hom(-,X), F)\cong F(X).$$
\end{lemma}
\begin{lemma}
\label{lemma-nat-x-y-hom-x-y}
	$\text{Nat}
	(\Hom(-,X),\Hom(-,Y))\cong\text{Hom}(X,Y)$.
\end{lemma}
\begin{remark}
\label{remark-corre-x-x-fully-faith-corre}
	The correspondence $X\mapsto\Hom(-,X)$ is
	fully faithful, that is, the correspondence
	$\Hom(X,X')\to\text{Nat}(\text{Hom}(-,X),
	\text{Hom}(-,X'))$
	is injective and bijective.
\end{remark}

\subsection{Brown's representability theorem}
\label{subsection-brown-s-repre-theor}

\subsubsection{In lecture13.pdf}
\label{subsection-in-lecture13-pdf}

See
\href{https://www.math.ru.nl/~gutierrez/%
files/Lecture13.pdf}{Lecture
13: Representable functors and the Brown
Representability Theorem.pdf}

Let us summarize the discussion so fat.
Suppose that
$$F:\text{Ho}(\text{Top}_*)^{\text{op}}\to
\text{Sets}_{*}$$
is a functor having the following two
properties:
\begin{enumerate}
	\item[(i)] $F\left( \bigvee_{i\in I}X_{i}
	\right) \to \prod_{i\in I}F(X_{i})$ is an
	isomorphism for any family of pointed spaces
	$\{X_{i}\} $.
\item[(ii)] "Pushout where one of the arrows
is a cofibration goes to weak pullback".
$F(B\cup_{A}C)\to F(B)\times_{F(A)} F(C)$ is
a surjection for any two cofibrations $A\to
B$ and $A\to C$.
\end{enumerate}

Then we also have
\begin{enumerate}
	\item[(iii)] $F(*)=*$
	\item[(iv)] If we have a pushout of pointed
	spaces
		$$
		\xymatrix{
			A\ar[r]\ar[d]&X\ar[d]\\
			B\ar[r]&Y
		}
		$$
		in which $i$ is a cofibration, then
		$F(Y)\to F(B)\times_{F(A)}F(X)$ is a
		sujrection.
	\item[(v)] If
	$\xymatrix@C=1.5em{A\ar@<.5ex>[r]\ar@<-
	.5ex>[r]&B\ar[r]&C}$
	is a coequalizer of pointed spaces in which
	the two maps form a cofibration $A\vee A\to
	B$, then $F(C)\to F(B)$ maps surjectively to
	the equalizer of
	$\xymatrix@C=1.5em{F(B)\ar@<.5ex>[r]\ar@<-
	.5ex>[r]&F(A)}$.
	\item[(vi)] If $X_{0}\to X_{1}\to X_{2}\to
	\ldots$ is a sequence of cofibrations, then
	$F(X)\to \lim_{\leftarrow}$ is a sujection.
\end{enumerate}

\begin{theorem}[Brown representability
theorem]
\label{theorem-brown-repre-theor}
	Let $F$ be a contravariant functor from the
	homotopy category of parallel connected
	CW-complexes to pointed sets. If $F$
	satisfies conditions (i) and (ii) above (for
	any pointed connected CW-complexes $X_{i}$,
	$A$, $B$, $C$), then $F$ is representable.
\end{theorem}

\begin{remark}[So what is a classifying
space?]
\label{remark-so-what-is-a-classifying-space}
	The theorem says that there is a space
	$B=B_{F}$ (itself a pointed CW-complex) for
	which there is a natural isomorphism
	$$\varphi :[X,B_{F}]_{*}\overset{\cong
	}{\longrightarrow}F(X)$$
	for any pointed CW-complex $X$. This space
	$B_{F}$ is called a {\it classifying space}
	for $F$. Recall also that when such
	$\varphi$ exists, it is completely
	determined by a \textit{generic} element
	$\gamma\in F(B_{F})$.

	The classifying space together with the
	genereic element is unique up to homotopy.
\end{remark}

Now let's sketch how the proof of Brown
representability theorem goes.

\begin{definition}
\label{definition-b-spher-class-space-f}
	$(B,\gamma)$ is a {\it spherical classifying
	space} for $F$ if $\gamma\in F(B)$ is an
	element inducing an isomorphism
	$$\gamma_{*}:[S^{n},B]_{*}\to
	F(S^{n}),\qquad \gamma_{*}(f)=f^{*}
	(\gamma)$$
	for any sphere $S^{n}$, $n>0$.
\end{definition}

\begin{remark}
\label{remark-if-b-1-1-b-2-2-two-such-sphe}
	If $(B_{1},\gamma_{1})$ and
	$(B_{2},\gamma_{2})$ are two such spherical
	classifying spaces in the category
	$\mathcal{C}$, and $f:B_{1}\to B_{2}$ is a
	map with $f^{*} (\gamma_{2})=\gamma_{1}$,
	then $f$ induces isomorphisms
	$$\pi_{n}(B_{1})\to \pi_{n}(B_{2})$$
	between the homotopy groups for each $n>0$.

	Since $B_{1}$ and $B_{2}$ are pointed
	connected CW-complexes, this means that $f$
	is a weak homotopy equivalence, hence a
	homotopy equivalence by Whitehead theorem.
\end{remark}

The proof of Brown representability theorem
is based in the following two lemmas:

\begin{lemma}
\label{lemma-let-x-be-point-cw-compl-let}
	Let $X$ be a pointed CW-complex and let $\xi
	\in F(X)$. Then there exists a spherical
	classifying space $(B,\gamma)$ for $F$ with
	a cofibration  $f:X\to B$ with $f^{*}
	(\gamma)=\xi$.
\end{lemma}

And in fact,

\begin{lemma}
\label{lemma-any-spher-class-space-b-f}
	Any spherical classifying space $(B,\gamma)$
	for $F$ is a classifying space.

	(Thus, $\gamma_{*}:[X,B]_{*}\to F(X)$ is an
	isomorphism for any pointed connected
	CW-complex $X$, not just for spheres.
\end{lemma}

\subsubsection{In Hatcher}
\label{subsection-in-hatcher}

\begin{theorem}[4.57]
\label{theorem-4-57}
	There are natural bijections
	$T:\left\langle{}X,K(G,n) \right\rangle{} \to H^{n}(X;G)$
	for all CW complexes $X$ and all $n>0$, with
	$G$ any abelian group. Such a $T$ has the
	form $T([f])=f^{*}$ for a certain
	distinguised class $\alpha\in
	H^{n}(K(G,n);G)$.
\end{theorem}

\begin{definition}[Fundamental class]
\label{definition-fundamental-class}
	A class $\alpha\in H^{n}(K(G,n);G)$ with the
	property stated in the theorem is called a
	{\it fundamental class}.
\end{definition}


\begin{theorem}[4E.1]
\label{theorem-4e-1}
	Every reduced cohomology theory on the
	category of basepointed CW complexes and
	basepoint-preserving maps has the form
	$h^{n} (X)=\left\langle{}X,K_{n} \right\rangle{} $ for some
	$\Omega$-spectrum $\{K_{n}\} $.
\end{theorem}


\subsection{Representability of generalized
cohomology theories}
\label{subsection-repre-gener-cohom-theor}

\subsection{$H^{n}(-,G)$ is represented by
$K(G,n)$ together with a chosen element in
$H^{n}(K(G,n),G)$}
\label{subsection-h-n-g-repre-k-g-n-toget}

Reduced cohomology is representable since

\begin{enumerate}
	\item[(i)]
	$\widetilde{H}^{n}(\bigvee_{i}X_{i})=
	\prod_{i}
	\widetilde{H}_{n}(X_{i})$.
	\item[(ii)] Also holds.
\end{enumerate}

What is the classifying space? Well, we have
$$[S^{p},B]=H^{n}(S^{p})=\begin{cases}
	0\qquad &p\neq n\\
	\mathbb{Z}\qquad &p=n
\end{cases}$$
which means that $B\cong K(\mathbb{Z},n)$.
More explicityly, we see that the homotopy
groups of the classifying space must coincide
with the reduced cohomology groups, which are
zero in all but the $n$-th dimension.

\subsection{Eilenberg-MacLane spaces can be
construced from Moore spaces by killing
higher homotopy groups}
\label{subsection-eilen-macla-space-can}
\subsection{Definition of $\Omega$-spectrum}
\label{subsection-definition-of-spectrum}


\subsection{Hurewicz homomorphism}
\label{subsection-hurewicz-homomorphism}
\begin{definition}[Hurewicz homomorphism]
\label{definition-hurewicz-homomorphism-2}
	$$
\begin{aligned}
		h: \pi_{n}(X) &\longrightarrow H_{n}(X) \\
		[f] &\longmapsto f_{*}(\alpha)
	.\end{aligned}
$$
\end{definition}
where $\alpha$ is a generator of $
H_{n}(S^{n})$.

\subsection{Absolute and relative versions of
Hurewicz theorem}
\label{subsection-absol-relat-versi-hurew}
\begin{theorem}[Hurewicz]
\label{theorem-hurewicz}
	If $X$ is $n$-connected for some $n\geq 2$,
	then $\tilde{H}_{i}(X)=0$ for $i<n$ and
	$\pi_{n}(X)\cong H_{n}(X)$.
\end{theorem}
\begin{lemma}
\label{lemma-if-x-y-simpl-conne-cw-space}
	If $X$ and $Y$ are simply connected
	CW-spaces and $f:X\to Y$ induces
	isomorphisms $H_{n}(X)\cong H_{n}(Y)$ for
	all $n$, then it is an homotopy equivalence.
\end{lemma}
\begin{theorem}[Relative Hurewicz]
\label{theorem-relative-hurewicz}
	If $(X,A)$ is $(n-1)$-connected, so
	$\pi_{n}(X,A)=0$ for $i\leq n-1$, for some
	$n\geq 2$ and $A$ is simply connected, then
	$H_{i}(X,A)=0$ for $i\leq n-1$ and
	$\pi_{n}(X,A)\cong H_{n}(X,A)$.

	\textit{Note. Perhaps $X$ should be
	$(n-1)$-connected for the relative version
	too---they are stated together in Hatcher.}
\end{theorem}
\begin{theorem}[Relative Hurewicz for general
$A$]
\label{theorem-relat-hurew-gener}
	If $(X,A)$ is an $(n-1)$-connected pair of
	path-connected spaces with $n\geq 2$ and
	$A\neq \emptyset $, then
	$h':\pi_{n}'(X,A,x_{0})\to H_{n}(X,A)$ is an
	isomorphism and $H_{i}(X,A)=0$ for $i<n$.
\end{theorem}



\section{Exact couple. Spectral sequence of
an exact couple.}
\label{section-exact-coupl-spect-seque}

These notes are based on S. Burkin lecture
notes and Hatcher.

\begin{enumerate}
	\item Suppose a space $X$ is expressed a
	union of a sequence of subspaces
		$$\ldots\subset X_p\subset X_{p+1}\subset
		\ldots X.$$
		Such a sequence is called a {\it filtration
		of $X$}. For example, $X$ could be a CW
		complex and $X_p$ its $p$-skeleton.

	\item Given such a filtration, for every $p$
	we have a pair $(X_p,X_{p-1})$ that yields a
	long exact sequence. All these long exact
	sequences can be arranged in a staricase
	fashion:
$$
\xymatrix@C=.7em@R=.8em{
	H_{p+q}(X_p)\ar@<.5ex>[r]\ar[r]\ar[d]_{i}
	&H_{p+q}(X_p/X_{p-1})\ar@<.5ex>[r]\ar[r]
	\ar[d]
	&H_{p+q-1}(X_{p-1})\ar[r]\ar[d]_{i}
	&H_{p+q-1}(X_{p-1}/X_{p-2})\ar[d]\ar[r]
	&\cdots\\
	H_{p+q}(X_{p+1})\ar[r]\ar[d]_{i}
	&H_{p+q}(X_{p+1}/X_{p})\ar[r]\ar[d]
	&H_{p+q-1}(X_{p})\ar@<.5ex>[r]\ar[r]_{j}
	\ar[d]_{i}
	&H_{p+q-1}(X_{p}/X_{p-1})\ar@<.5ex>[r]\ar[d]
	\ar[r]_{k}
	&\cdots\\
	\cdots&\cdots&\cdots&\cdots
}
$$
	\item If you put all the absolute groups
	$H_n(X_p)$ in a big $A=\bigoplus_{n,p\in
	\mathbb{N}}H_n(X_p)$ and all the relative
	$H_n(X_p,X_{p-1})$ in a big
	$E=\bigoplus_{n,p\in
	\mathbb{N}}H_n(X_p,X_{p-1})$, you can write
	all those long exact sequences in a compact
	triangle diagram called an {\it exact
	couple} triangle diagram
	$$
	\xymatrix{
		A\ar[rr]^{i}&&A\ar[dl]^{j}\\
		&E\ar[ul]^{k}
	}
$$
here the maps $i,j$ and $k$ are just the maps
from the long exact sequences in the other
big diagram (which makes the exact couple
actually \textit{exact}).

	\item There is a nice way to define new
	groups $A'$ and $E'$ and maps between them
	that yield a {\it derived couple}, which is
	another exact couple.	Namely,
		\begin{itemize}
			\item $A^1_{n,p}:=H_n(X_p)$
			\item $E^1_{n,p}:=H_n(X_p,X_{p-1})$.
			\item $d:=jk$
			\item $E':=\ker d/\text{img}d$, the
			homology of $E$ respect to $d$.
			\item $A':=i(A)\subset A$.
			\item $i':=i|_{A'}$.
			\item $j'(ia):=[ja]\in E'$.
			\item $k'[e]:=ke$.
		\end{itemize}
	anyway, they are well defined and yield the
	exact
	$$
	\xymatrix@C=1.5em@R=1em{
		A'\ar[rr]^{i'}&&A'\ar[dl]^{j'}\\
		&E'\ar[ul]^{k'}
	}
	$$

	\item The sequence $E,E',\ldots$ is called a
	{\it spectral sequence}.
\end{enumerate}

\section{First quadrant. Convergence}
\label{section-first-quadrant-convergence}

[Explanation that when the spectral sequence
is defined only in the first quadrant it
eventually stabilizes and $E_{\infty}$-page
is well-defined.]

\section{Homological and cohomological
spectral sequences}
\label{section-homol-cohom-spect-seque}

\noindent
\emph{[The introductory discussion for this section is missing from
the notes.]}

\begin{remark}
\label{remark-diffe-spect-seque-homol-go}
	The differentials in the spectral sequence
	for homology go up-left and for cohomology
	they go right-down.
\end{remark}

\section{Serre spectral sequence}
\label{section-serre-spectral-sequence}

\section{Homological version}
\label{section-homological-version}

\begin{theorem}[Serre spectral sequence for
homology]
\label{theorem-serre-spect-seque-homol}
	Let $F\to X\to B$ be a fibration with $B$
	path-connected. If $\pi_1(B)$ acts trivially
	on $H_*(F;G)$, then there is a spectral
	sequence $\{E^r_{p,q},d_r\}$ with
	\begin{itemize}
		\item
			$$d^r_{p,q}:E^r_{p,q}\to
			E^r_{p-r,q+r-1}\qquad \text{and} \qquad
			E^{r+1}_{p,q}=\ker d^r_{p,q}/\text{img
			}d^r_{?}$$
		\item The groups
			$$F_pH_n:=\text{img}(H_n(X_p)\to H_n(X))$$
		where the map $H_{n}(X_p)\to H_{n}(X)$ is
		just the induced map by inclusion
		$X_p\hookrightarrow X$, form a filtration
			% \item Stable terms $E^\infty_{p,n-p}$
			% are isomorphic to the successive
			% quotients $F_n^p/F^{p-1}_n$ in a
			% filtration
		$$0\subset F_0H_n\subset \ldots\subset
		F_nH_n=H_n(X;G)$$
		of $H_n(X;G)$ such that
		$$E^\infty_{p,n-p}\cong
		F_pH_n/F_{p-1}H_n.$$
		Another way to write this is
		$$E_{p,q}^\infty=F_pH_{p+q}\Big/F_{p-1}
		H_{p+q}.$$
	\item $$E^2_{p,q}\cong H_p(B;H_q(F;G))$$
\end{itemize}
\end{theorem}

\begin{remark}
\label{remark-case-serre-fibra-p-n-map}
	In the case of a Serre fibration, for $p\gg
	n$ the map $H_{n}(X_p)\to H_{n}(X)$ is an
	isomorphism. The group
	$\bigoplus_{p}F^p_n/F^{p-1}_n$ is called the
	{\it associated graded group of the filtered
	group $H_n(X)$}.
\end{remark}

\begin{remark}[Quote from Hatcher]
\label{remark-quote-from-hatcher}
	If the coefficient group $G$ is a field then
	$H_{n}(X;G)$ is the direct sum
	$\bigoplus_{p} E^\infty_{p,n-p}$.
\end{remark}

\begin{remark}
\label{remark-can-we-recov-h-n-e-gener}
	Can we recover $H^{n}(E)$ in general? No.
\end{remark}

\section{Cohomological version}
\label{section-cohomological-version}

\begin{theorem}[Serre spectral sequence for
cohomology]
\label{theorem-serre-spect-seque-cohom}
	Let $F\to X\to B$ be a Serre fibration with
	$\pi_0(B)=0$ such that the action of
	$\pi_1(B)$ on $H^*(F;G)$ is trivial.

	Then there is a spectral sequence
	$\{E_r^{p,q},d_r\}$ such that
\begin{itemize}
	\item $$d_r:E_r^{p,q}\to
	E_r^{p+r,q-r+1}\qquad\text{and} \qquad
	E_{r+1}^{p,q}=\ker d_r/\text{img }d_r.$$
	\item The groups
		$$F^pH^n:=\text{img} (F^pH^n\to
		F^nH^n=H^n(X))$$
	form a filtration
	$$H^{n}(X)=F^0H^n\supset F^1H^n\supset
	\ldots\supset F^nH^n\supset \emptyset $$

	of $H^{n}(X,G)$ such that
	$$E_{\infty}^{p,q}\cong
	F^pH^{p+q}/F^{p+1}H^{p+q}.$$
\item $$E^{p,q}_2\cong H^p(B,H^{q}(F,G)).$$
\end{itemize}
\end{theorem}

\section{Product in cohomological Serre
spectral sequence}
\label{section-produ-cohom-serre-spect}
Hatcher, Spectral sequences p. 543.

The Serre spectral sequence for cohomology
becomes much more powerful when cup products
are brought into the picture. For this we
need to consider cohomology with coefficients
in a ring $R$ rather than just a group $G$.
What we will show is that the spectral
sequence can be provided with bilinear
products
$$E^{p,q}_{r}\times E^{s,t}_{r}\to
E^{p+s,q+t}_{r}$$
for $1\leq r\leq \infty$ satisfying the
following properties
\begin{itemize}
	\item Each differential $d_{r}$ is a
	derivation, satisfying
		$$d(xy)=(dx)y+(-1)^{p+q} x(dy)\qquad
		\forall x\in E^{p,q}_{r}$$
	This implies that the product on page $r$
	induces a product in page $r+1$, and the
	product in page $\infty$ shall be the one
	induced by previous products.
	\item The product on the second page,
	$E^{p,q}_{2}\times E^{s,t}_{2}\to
	E^{p+s,q+t}_{2}$ is $(-1)^{qs}$ times the
	standard (standard? because as I remember
	cup product did not move the coeffients,
	right?) cup product
		$$H^{p}(B;H^{q}(F;R))\times
		H^{s}(B;H^{t}(F;R))\to
		H^{p+s}(B;H^{q+t}(F;R))$$
		sending a pair of cocycles $(\phi,\psi)$ to
		$\phi\smile \psi$ where coefficients are
		multiplied via the cup product
		$H^{q}(F;R)\times H^{t}(F;E)\to
		H^{q+t}(F;R)$.
	\item The cup product in $H^{*}(X;R)$
	restricts to maps $F^{m}_{p}\times
	F^{n}_{s}\to F^{m+n}_{p+s}$. These induce
	quotient maps $F^{ m}/F^{m}_{p+1}\times
	F^{n}_s/F^{n}_{p+s}\to
	F^{m+n}_{p+s}/F^{m+n}_{p+s+1}$ that coincide
	with the product $E^{p,m-p}_{\infty}\times
	E^{s,n-s}_{\infty}\to
	E^{p+s,m+n-p-s}_{\infty}$.
\end{itemize}

\section{Wang spectral sequence}
\label{section-wang-spectral-sequence}

Let's use lecture notes and
\href{https://en.wikipedia.org/wiki/%
Spectral_sequence#Wang_sequence}{wikipedia}
to understand Wang spectral sequence.

Consider a fibration over a sphere
$F\hookrightarrow E\to B$. According to our
theorem for Serre spectral sequences, we know
that the $n$-th page looks as follows:
$$
\begin{array}{c|cccccccc}
	n&H_n(F)&&&&&&&H_{n}(F)\\
	 &\vdots&&&&&&&\vdots\\
	2&H_2(F)&&&&&&&H_2(F)\\
	1&H_1(F)&&&&&&&H_1(F)\\
	0&H_0(F)&&&&&&&H_0(F)\\
	\hline\\
	 &0&&&\ldots&&&&n
\end{array}
$$

We can see there can be nontrivial
differentials only on the $n$th page, so that
$E^{n+1}=E^n$. Moreover, the differentials
that don't vanish is are of the form
$$
\xymatrix{
E^n_{n,q}\ar[r]^{d^n_{n,q}}&E^n_{0,q+n-1}
}
$$

For $d_{n,q-n}^n$ we obtain
$$
\xymatrix{
	0\ar[r]&\ker
	d^n_{n,q-n}\ar@{^{(}->}[r]^{\text{inclusion}
	}
	&E^n_{n,q-n}\ar[r]^{d^n_{n,q-n}}
	&E^n_{0,q-1}\ar@{->>}[r]^{\text{quotient}}
	&\text{coker}d^n_{n,q-n}\ar[r]&0
}
$$
But all that is just
$$
\xymatrix@R=1em{
	0\ar[r]&\ker
	d^n_{n,q-n}\ar@{^{(}->}[r]^{\text{inclusion}
	}
	\ar[d]^{=}
	&E^n_{n,q-n}\ar[r]^{d^n_{n,q-n}}\ar[d]^{=}
	&E^n_{0,q-1}\ar@{->>}[r]^{\text{quotient}}
	\ar[d]^{=}
	&\text{coker}d^n_{q-n}\ar[r]\ar[d]^{=}&0\\
	0\ar[r]&\dfrac{\ker
	d^n_{n,q-n}}{\underbrace{\text{img
	}d^{n-1}_{n,q-n}}_{=0}}
	\ar[r]\ar[d]^{=}
	&E^2_{n,q-n}\ar[r]\ar[d]^{=}
	&E^2_{0,q-1}\ar[r]\ar[d]^{=}
	&\dfrac{E^n_{n,q-n}}{\text{img }d^n_{n,q-n}}
	\ar[d]^{=}\ar[r]&0\\
	0\ar[r]&E^\infty_{n,q-n}\ar[r]&H_{q-n}(F)
	\ar[r]
	&H_{q-1}(F)\ar[r]&E^{\infty}_{0,q-1}\ar[r]&0
}
$$
This is the first "half" of the Wang
sequence. For the other half recall that by
the theorem of Serre spectral sequence for
homology we know that
$$E^{\infty}_{p,n-p}=F_pH_n/F_{p-1}H_n$$
for a filtration on the $n$-th homology group
of the total space
$$\emptyset=F_{-1}\subset F_0 \subset
F_1\subset \ldots\subset F_n=H_{n}(E).$$
And recall that we may also write
$$E^\infty_{p,q}=F_pH_{p+q}/F_{p-1}H_{p+q}.$$
Then we can see that
$$
\xymatrix@R=1em{
	0\ar[r]&E^\infty_{0,q}\ar[d]^{=}\ar[r]^{i^*}
	&H_{q}(E)\ar[r]&E^\infty_{n,q-n}\ar[r]\ar[d]
	^{=}
	&0\\
	&\dfrac{F_0H_{q}(F)}{\underbrace{F_{-1}H_{q}
	(F)}_{=0}}
	&&\dfrac{F_nH_{q}(E)}{F_{n-1}H_{q}(E)}
}
$$
But
$$E^\infty_{0,q}=\dfrac{F_0H_{q}(F)}
{\underbrace{F_{-1}H_{q}(F)}_{=0}}\quad
\overset{\text{why?}}{=}H_{q}(F)$$

This is the other "half" of the Wang
sequence. It only remains to put both halves
together and conclude that
$$
\xymatrix@C=1.5em{
	\cdots\ar[r]&H_{q}(F)\ar[r]^{i^*}&H_{q}(E)
	\ar[r]
	&E^{\infty}_{n,q-n}\ar[r]^{d^n_{n,q-n}}
	&H_{q-n}(F)\ar[r]&H_{q-1}(F)\ar[r]
	&E^\infty_{0,q-1}\ar[r]&\cdots
}
$$
But we may remove the $E^\infty$ terms to get
$$
\xymatrix@C=1.5em{
	\cdots\ar[r]&H_{q}(F)\ar[r]^{i^*}&H_{q}(E)
	\ar[r]
	&H_{q-n}(F)\ar[r]^{d^n_{n,q-n}}&H_{q-1}(F)
	\ar[r]
	&H_{q-1}(E)\ar[r]&H_{q-n-1}(F)\ar[r]&\cdots
}
$$

\section{Extra: spectral sequences for chain
complexes and the Fr\"olicher spectral
sequence}
\label{section-extra-spect-seque-chain}
Let's follow Griffiths and Harris to
construct a nice generalization of this
theory.

A {\it chain complex} $(K^\bullet,d)$ is a
sequence of abelian groups (if working in the
category $\text{ab}$)
$$
\xymatrix@C=1.5em{
	K^{0}\ar[r]^{d}&K^{1}\ar[r]^{d}&K^2\ar[r]
	^{d}&\cdots
}
$$
such that the differential $d$ satisfies
$d\circ d=0$. We have {\it cohomology groups}
$$H^{p}(K^*)=\frac{\ker d:K^p\to
K^{p+1}}{\text{img }d:K^{p-1}\to K^p}$$
and the {\it cohomology ring}
$$H^{*}(K^*)=\bigoplus_{p\geq 0}
H^{p}(K^*).$$

A {\it subcomplex} $(J^*,d)$ is given by
subgroups $J^*\subset K^*$
 with $dJ^*\subset J^*$.
The {\it quotient complex} $(L^*,d)$ is also
well-defined by $L^*=K^*/J^*$ and the obvious
differential. this construction yields a {\it
short exact sequence of complexes}
$$
\xymatrix{
	0\ar[r]&J^*\ar[r]&K^*\ar[r]&L^*\ar[r]&0
}
$$
and by the fundamental theorem of homological
algebra, this yields
$$
\xymatrix{
	\cdots\ar[r]&H^{p}(J^*)\ar[r]&H^{p}(K^*)
	\ar[r]
	&H^{p}(K^*)\ar[r]&H^{p+1}(J^*)\ar[r]&\cdots
}
$$

Generalizing the notion of subcomplex, a {\it
filtered complex} $(F^pK^*,d)$ is a
decreasing sequence of subcomplexes
$$K^*=F^0K^*\supset
F^1K^*\supset\cdots\supset F^nK^*\supset
F^{n+1}K^*=\{0\}$$
and the spectral sequence will generalize the
long exact sequence in cohomology.

The {\it associated graded complex} to a
filtered complex $ (F^pK^*)$ is the complex
$$\text{Gr}K^*=\bigoplus_{p>0}
F^pK^*/F^{p+1}K^* $$
with the obvious differential. The filtration
$F^pK^*$ on $K^*$ also induces a filtration
$F^pH^*(K^*)$ on cohomology by (I think)
$$F^pH^q(K^*)=\frac{\ker d:F^pK^q\to
F^pK^{q+1}}{\text{img }d:F^pK^{q-1}\to
F^pK^q}$$
This means that the $F^pH^q(K^*)$ are
subgroups of $H^{q}(K^*)$. This filtration in
turn induces an {\it associated graded in
cohomology}
$$\text{Gr}H^*(K^*)=\bigoplus_{p,q}
F^pH^q(K^*)\big/F^{p+1}H^q(K^*)
$$

\begin{definition}
\label{definition-spect-seque-seque-e-r}
	A {\it spectral sequence} is a sequence
	$\{E_r,d_r\}$ with $r\geq 0$ of bigraded
	groups
	$$E_r=\bigoplus_{p,q\geq 0} E^{p,q}_r $$
together with differentials
$$d_r:E^{p,q}_r\to E^{p+r,q-r+1}_r,\qquad
d^2_r=0$$
such that
$$H^{*}(E_r)=E_{r+1}$$
I just wonder what is the cohomology of a
bigraded group, right? there are several
chain complexes and cohomologies, but ok…
\end{definition}

\begin{proposition}[existence and description
of the chain complex spectral sequence]
\label{proposition-exist-descr-chain-compl}
	Let $K^*$ be a filtered complex. then there
	exists a spectral sequence $\{E_r\}$ with
	$$
\begin{aligned}
		E^{p,q}_0&=F^pK^{p+q}/F^{p+1}K^{p+q},\\
		E^{p,q}_1&=H^{p+q}\left( F^pK^*/F^{p+1}K^*
		\right)\\
		E^{p,q}_\infty&=F^pH^{p+q}(K^*)\big/F^{p+1}
		H^{p+q}(K^*)
	\end{aligned}
$$
and we just say that the spectral sequence
{\it abuts} to $H^{*}(K^*)$.
\end{proposition}

\paragraph{Explanation.} We have defined four
objects:
\begin{enumerate}
	\item A filtered complex.
	\item A graded complex associated to the
	filtered complex whose entries are direct
	sums of quotients of the filtered complex.
	\item A filtration on the cohomology of the
	complex given by the cohomologies of every
	complex of the filtered complex.
	\item A graded complex associated to the
	filtration on the cohomology whose entries
	are direct sums running over the cohomology
	groups and the filtration entries.
\end{enumerate}
And the proposition says is that there is a
spectral sequence such that
\begin{enumerate}
	\item The $(p,q)$-entry of the zero page is
	the (first) graded at $p+q$.
	\item The $(p,q)$-entry of the first page is
	the $p+q$ cohomology of the graded.
	\item The $(p,q)$-entry of the $\infty$-page
	is the graded of the cohomology
 (ie. the second graded).
\end{enumerate}

Now consider the following situation. A {\it
double complex} is a bigraded group
$$K^{*,*}=\bigoplus_{p,q\geq 0}  K^{p,q}$$
equipped with two differentials
$$d:K^{p,q}\to K^{p+1,q}$$
$$\delta:K^{p,q}\to K^{p,q+1}$$
satisfying $d^2=\delta^2=0$ and $d\delta
+\delta d=0$.

There is an {\it associated single complex}
$(K^*,d)$ defined by
$$K^n=\bigoplus_{p+q=n} K^{p,q} $$
with differential
$$D=d+\delta.$$
there are two filtrations on the associated
single complex:
$$'F^pK^n=\bigoplus_{\substack{p'+q=n\\p'\geq
p}}K^{p',q},$$
	$$''F^qK^n=\bigoplus_{\substack{p+q''=
	n\\q''\geq
	q}}K^{p,q''}.$$
each of these filtrations yields a spectral
sequence according to our proposition, thus
both abutting to the cohomology of the total
space $H^{*}(K^*)$.

Now we turn to Voisin:

\begin{proposition}[Voisin, 8.25]
\label{proposition-voisin-8-25}
	Let $K^\cdot$ be the simple complex
	associated to a double complex
	$(K^{p,q},D_1,D_2)$. Then if
	$F^pK^n=\bigoplus_{r\geq o,r+s=n} K^{r,s}$
	is the filtration introduced in section
	8.3.1, the spectral sequence associated to
	$F$ has first terms given by
	\begin{itemize}
		\item $E^{p,q}_0=K^{p,q}$, $d_0=(-1)^pD_2$.
		\item $E^{p,q}_1=H^{q}(K^{p,\cdot})$,
	\end{itemize}
	and the differential $d_1:H^{q}(K^{p,\cdot
	})\to H^{q}(K^{p+1})$ is induced by the
	morphism of complexes $D_1:K^p\to K^{p+1}$.

\end{proposition}

Now to the case that concerns us. Let's go
back to Griffiths and Harris. Let $M$ be a
complex manifold and define double complex
$$K^{p,q}=A^{p,q}(M)$$
with differentials $d$ and $\bar\partial$. I
wonder why the first one is $ d$… it appears
more intuitive to choose $\partial$.

The associated single complex is just the de
Rham complex since
$$A^{n}=\bigoplus_{p+q=n} A^{p,q}.$$
The de Rham differential $d$ satisfies
$d=\partial+\bar\partial$, so I definitely
would like to take $\partial$ instead of $d$
as I said above… I guess I'm not certain
whether $\partial$ induces a chain complex, or
if
$\partial\bar\partial+\bar\partial\partial$.
Then the corresponding filtrations yield the
{\it Fr\"olicher spectral sequences}, both of
which abut to $H^{*}_{\text{dR}}(m)$.

Now if $M$ is compact K\"ahler then every
class $[a]\in 'E^{p,q}_1\cong
H^{p,q}_{\bar\partial}(m)$ has a harmonic
representative for the
$\bar\partial$-laplacian (following
Huybrechts,
$\Delta_{\bar\partial}=
\bar\partial^*\bar\partial+
\bar\partial\bar\partial^*$).
Shouldn't it be the \textit{second}
filtration, $''E^{p,q}$, ie. the one
associated to $\bar\partial$?

Now recall one the K\"ahler identities (prop.
3.1.12 in Huybrechts):
$$2\Delta_{\bar\partial}=\Delta_d:=d^*d+dd^*,
$$
which implies that $da=0$. I would like to
check why this holds, ie. why
$\bar\partial$-harmonic $\implies$
$d$-harmonic $\implies$ closed. This means
that the differentials of the spectral
sequence associated to the first
differential, $d$, are all zero, so
\textit{the spectral sequence degenerates in
the first page}, that is
$$'E_1\cong 'E_2\cong \ldots\cong
'E_\infty.$$
in this case the corresponding filtration is
the {\it Hodge filtration}:
$$
\begin{aligned}
	F^pH^{n}_{\text{dR}}(M)&\cong
	H^{n,0}(m)\oplus\ldots\oplus H^{p,n-p}(M)\\
				       &=\bigoplus_{\substack{p'+q=n \\
				       p'\geq p}}  H^{p',q}(M)
\end{aligned}
$$
\begin{remark}
\label{remark-if-m-not-k-ahler-spect-seque}
If $M$ is not K\"ahler the spectral sequence
need not degenerate at the first page. the
{\it Iwasawa manifold} given by a quotient of
a Lie group of complex matrices is an example
where $'E_1\neq 'E_2$, but no example is
known there ${'E_2\neq 'E_\infty}$
\end{remark}

\begin{remark}[Altan]
\label{remark-altan}
	It appears that we have only argued that the
	differential on the first page of the
	Fr\"olicher spectral sequence is zero. Other
	spectral seqeunces like the Gysin spectral
	sequence also have vanishing differentials
	on the first pages, but this does not imply
	that the first and the infinity page
	coincide. How can we explain that this is
	actually the case for the Fr\"olicher
	spectral sequence?
\end{remark}

\paragraph{Answer}It looks to me like the
proof reduces to Dolbeault´s theorem. Not
only on the first page can we use that any
form has a $\bar\partial$-harmonic
representative so that the differential
$d_1=\bar\partial$ (or $\partial$?) vanishes,
but also on any other page since all the
$E_{r}$ are just subquotients of some
Dolbeault cohomology group---in the end we
just gave take a representative (of the
representative of the representative…) that
is $\bar\partial$-harmonic. And this all just
made me think:

\begin{remark}[Dani]
\label{remark-dani}
	Then what was all that
	$\bar\partial$-harmonic implies $d$-closed
	discussion in Griffiths and Harris for?
\end{remark}

\paragraph{Atlan's answer to Altan's
question} We have shown that
$$E_1=\bigoplus_{p,q} H^{q}(X,\Omega^q)$$
and we also know that
$$E_\infty=\bigoplus_{p,q} E_\infty^{p,q} $$
So we get an isorphism between these pages by
the fact that the filtration on the de Rham
complex is the Hodge filtration. Voisin
justifies this via
\begin{proposition}[Voisin, 7.5]
\label{proposition-voisin-7-5}
Let $F^p\mathcal{A}^k(X)$ be the set of
complex differential forms which are sums of
forms of type $(r,k-r)$ with $r\geq p$ at
every point. Then
$$F^pH^{k}(X,\mathbb{C})=\dfrac{\ker
d:F^pA^k(X)\to F^pA^{k+1}(X) }{\text{img
}d:F^pA^{k-1}(X)\to F^pA^k(X)}$$
\end{proposition}

\begin{remark}
\label{remark-why-all-meani-voisi-claim}
	Why is all of this meaningful? Voisin claims
	that this result is very nearly
	equivalent to the Hodge theorem, which says
	that
\begin{equation*}\label{eq:hodge-decom}
	 H^{k}(X,\mathbb{C})=\bigoplus_{p+q=k}
	 H^{p,q}(X).
\end{equation*}
While this precise equation cannot be derived
from the Fr\"olicher spectral sequence, what
we \textit{did} obtain is
$$\bigoplus_{p,q}F^pH^{p+q}(X,\mathbb{C})/
F^{p+1}H^{p+q}(X,\mathbb{C})=E^{p,q}_{\infty}
=E^{p,q}_{1}=H^q(X,\Omega^p).$$

\paragraph{Question.}
Why are Hodge theorem and the conclusion from
the degeneration of the Fr\"olicher spectral
sequence nearly equivalent?
\end{remark}





\section{Cohomology of spaces}
\label{at-section-cohomology-spaces}

\noindent
Let $X$ be a topological space and $G$ an
abelian group.  The singular chain complex
is
$$
\xymatrix{
C_\bullet(X)
  & \cdots \ar[r]
  & C_{n+1}(X)
    \ar[r]^-{\partial_{n+1}}
  & C_n(X) \ar[r]^-{\partial_n}
  & C_{n-1}(X) \ar[r]
  & \cdots.
}
$$
Here $C_n(X)$ is the free abelian group on
the singular simplices
$\sigma:\Delta^n\to X$.  Thus
$$
C_n(X)
=
\bigoplus_{\sigma:\Delta^n\to X}
\mathbb{Z}\sigma.
$$
We define
$$
H_n(X;\mathbb{Z})=H_n(C_\bullet(X)),
$$
$$
C^n(X;G)
=
\operatorname{Hom}(C_n(X),G),
$$
and
$$
H^n(X;G)=H^n(C^\bullet(X;G)).
$$
The cochain complex runs in the opposite
direction:
$$
\xymatrix{
\cdots
  & \operatorname{Hom}(C_{n+1},G)
    \ar[l]
  & \operatorname{Hom}(C_n,G)
    \ar[l]_-{\delta^n}
  & \operatorname{Hom}(C_{n-1},G)
    \ar[l]_-{\delta^{n-1}}
  & \cdots \ar[l].
}
$$
The coboundary of a cochain is given by
the diagonal composite
$$
\xymatrix{
C_n(X) \ar[r]^-\varphi
  & G \\
C_{n+1}(X)
  \ar[u]^{\partial_{n+1}}
  \ar[ur]_{\delta\varphi}.
}
$$
Elements of $\ker(\delta)$ are
{\it cocycles}; elements of
$\operatorname{im}(\delta)$ are
{\it coboundaries}.

\noindent
The universal coefficient theorem gives
$$
\xymatrix@C=.55em{
0 \ar[r]
  & \operatorname{Ext}^1
    (H_{n-1}(X),G) \ar[r]
  & H^n(X;G) \ar[r]^-h
  & \operatorname{Hom}(H_n(X),G)
    \ar[r]
  & 0.
}
$$

\noindent
Reduced cohomology is obtained by dualizing
the augmented singular complex
$$
\xymatrix@C=.6em{
\cdots \ar[r]
  & C_2(X) \ar[r]^-{\partial_2}
  & C_1(X) \ar[r]^-{\partial_1}
  & C_0(X) \ar[r]^-\varepsilon
  & \mathbb{Z} \ar[r]
  & 0.
}
$$
For $n>0$ there is an isomorphism
$$
\widetilde H^n(X;G)\cong H^n(X;G).
$$
In degree zero,
$$
H^0(X;G)
\cong
\prod_{[x]\in\pi_0(X)}G.
$$
If $X$ is path connected, the Hurewicz map
and the universal property of abelianization
give
$$
H^1(X;G)
\cong
\operatorname{Hom}(\pi_1(X),G).
$$
There is also a reduced universal
coefficient theorem.

\noindent
For a pair $A\subset X$, begin with the
split short exact sequence
$$
\xymatrix{
0 \ar[r]
  & C_\bullet(A) \ar[r]
  & C_\bullet(X) \ar[r]
  & C_\bullet(X,A) \ar[r]
  & 0.
}
$$
It splits degree by degree because the
singular simplices in $A$ form a subset of
the basis of $C_\bullet(X)$.  Dualizing
gives
$$
\xymatrix{
0 \ar[r]
  & C^\bullet(X,A;G) \ar[r]
  & C^\bullet(X;G) \ar[r]
  & C^\bullet(A;G) \ar[r]
  & 0.
}
$$
We define
$$
H^n(X,A;G)
=
H^n(C^\bullet(X,A;G)).
$$
The long exact sequence of the pair is
$$
\xymatrix@C=.55em{
\cdots \ar[r]
  & H^n(X,A;G) \ar[r]
  & H^n(X;G) \ar[r]
  & H^n(A;G) \ar[r]
  & H^{n+1}(X,A;G) \ar[r]
  & \cdots.
}
$$

\noindent
Let $f:(X,A)\to(Y,B)$ be a map of pairs.
The induced chain maps form a commutative
diagram.
$$
\xymatrix@C=.65em{
0 \ar[r]
  & C_\bullet(A) \ar[r] \ar[d]_{f_\#}
  & C_\bullet(X) \ar[r] \ar[d]_{f_\#}
  & C_\bullet(X,A) \ar[r]
    \ar[d]_{f_\#}
  & 0 \\
0 \ar[r]
  & C_\bullet(B) \ar[r]
  & C_\bullet(Y) \ar[r]
  & C_\bullet(Y,B) \ar[r]
  & 0.
}
$$
Dualizing reverses the arrows:
$$
\xymatrix@C=.65em{
0 \ar[r]
  & C^\bullet(X,A;G) \ar[r]
  & C^\bullet(X;G) \ar[r]
  & C^\bullet(A;G) \ar[r]
  & 0 \\
0 \ar[r]
  & C^\bullet(Y,B;G)
    \ar[r] \ar[u]^{f^\#}
  & C^\bullet(Y;G)
    \ar[r] \ar[u]^{f^\#}
  & C^\bullet(B;G)
    \ar[r] \ar[u]^{f^\#}
  & 0.
}
$$
The long exact sequences are natural:
$$
\xymatrix@C=.45em{
\cdots \ar[r]
  & H^n(X,A;G) \ar[r]
  & H^n(X;G) \ar[r]
  & H^n(A;G) \ar[r]
  & H^{n+1}(X,A;G) \ar[r]
  & \cdots \\
\cdots \ar[r]
  & H^n(Y,B;G) \ar[r] \ar[u]^{f^*}
  & H^n(Y;G) \ar[r] \ar[u]^{f^*}
  & H^n(B;G) \ar[r] \ar[u]^{f^*}
  & H^{n+1}(Y,B;G)
    \ar[r] \ar[u]^{f^*}
  & \cdots.
}
$$
In particular,
$$
f^*:H^n(Y,B;G)\longrightarrow H^n(X,A;G)
$$
is contravariantly functorial.

\noindent
Cohomology is homotopy invariant.  If two
maps of pairs are homotopic, the induced
chain maps are chain homotopic.  Dualizing
the chain homotopy gives a cochain
homotopy, so the induced cohomology maps
are equal.

\noindent
The universal coefficient theorem for a
pair reads
$$
\xymatrix@C=.55em{
0 \ar[r]
  & \operatorname{Ext}^1
    (H_{n-1}(X,A),G) \ar[r]
  & H^n(X,A;G) \ar[r]
  & \operatorname{Hom}(H_n(X,A),G)
    \ar[r]
  & 0.
}
$$
It applies because $C_\bullet(X,A)$ is a
complex of free abelian groups.  Its
naturality diagram is
$$
\xymatrix@C=.3em@R=2.3em{
0 \ar[r]
  & \operatorname{Ext}^1
    (H_{n-1}(X,A),G) \ar[r]
  & H^n(X,A;G) \ar[r]^-h
  & \operatorname{Hom}(H_n(X,A),G)
    \ar[r]
  & 0 \\
0 \ar[r]
  & \operatorname{Ext}^1
    (H_{n-1}(Y,B),G)
    \ar[r] \ar[u]
  & H^n(Y,B;G)
    \ar[r]_-{h'} \ar[u]^{f^*}
  & \operatorname{Hom}(H_n(Y,B),G)
    \ar[r] \ar[u]
  & 0.
}
$$

\noindent
Cohomological excision follows from
homological excision and the five lemma.
If $Z\subset A\subset X$ and
$\overline Z\subset\operatorname{int}(A)$,
then inclusion induces an isomorphism
$$
H^n(X,A;G)
\longrightarrow
H^n(X-Z,A-Z;G).
$$

\noindent
For a wedge whose gluing point has a
contractible neighborhood,
$$
\widetilde H_n
\left(\bigvee_\alpha X_\alpha\right)
\cong
\bigoplus_\alpha
\widetilde H_n(X_\alpha)
$$
and
$$
\widetilde H^n
\left(\bigvee_\alpha X_\alpha;G\right)
\cong
\prod_\alpha
\widetilde H^n(X_\alpha;G).
$$
For disjoint unions,
$$
H_n\left(\coprod_\alpha X_\alpha\right)
\cong
\bigoplus_\alpha H_n(X_\alpha)
$$
and
$$
H^n
\left(\coprod_\alpha X_\alpha;G\right)
\cong
\prod_\alpha H^n(X_\alpha;G).
$$
\section{Cellular cohomology}
\label{at-section-cellular-cohomology}

\noindent
For a CW complex $X$,
let $I_n$ denote the set of its $n$-cells.
$$
H_n(X^n,X^{n-1})
\cong
\bigoplus_{\alpha\in I_n}\mathbb{Z}.
$$
The cellular boundary
is the composite shown by the upper
diagonals in
{\small
$$
\xymatrix@C=0em@R=1.8em{
  & H_{n-1}(X^{n-1})
    \ar[rd]^{j_{n-1}} \\
H_n(X^n,X^{n-1})
  \ar[ru]^{\partial_n}
  \ar[rr]^-{d_n}
  & & H_{n-1}(X^{n-1},X^{n-2})
      \ar[rd]_{\partial_{n-1}}
      \ar[rr]^-{d_{n-1}}
  & & H_{n-2}(X^{n-2},X^{n-3}) \\
  & & & H_{n-2}(X^{n-2})
          \ar[ru]_{j_{n-2}}.
}
$$
}
The compositions $d_{n-1}d_n$ vanish by
the long exact sequences of the relevant
triples.  The individual cellular boundary
maps need not vanish.

\noindent
Dualizing gives the cellular cochain
differentials.
{\small
$$
\xymatrix@C=0em@R=1.8em{
  & H^{n-1}(X^{n-1})
    \ar[ld]_{\partial_n^*} \\
H^n(X^n,X^{n-1})
  & & H^{n-1}(X^{n-1},X^{n-2})
      \ar[ll]_-{D^{n-1}}
      \ar[lu]_{j_{n-1}^*}
  & & H^{n-2}(X^{n-2},X^{n-3})
      \ar[ll]_-{D^{n-2}}
      \ar[ld]^{j_{n-2}^*} \\
  & & & H^{n-2}(X^{n-2})
          \ar[lu]^{\partial_{n-1}^*}.
}
$$
}

\begin{proposition}
\label{at-proposition-cellular-cohomology}
Let $X$ be a CW complex.  The cellular
cochain complex
$$
C^n_{\mathrm{CW}}(X;G)
=
\operatorname{Hom}
(C_n^{\mathrm{CW}}(X),G)
$$
computes singular cohomology:
$$
H^n(C^\bullet_{\mathrm{CW}}(X;G))
\cong H^n(X;G).
$$
\end{proposition}

\begin{proof}
Since
$$
C_n^{\mathrm{CW}}(X)
\cong
\bigoplus_{\alpha\in I_n}\mathbb{Z},
$$
we have
$$
\operatorname{Hom}
(C_n^{\mathrm{CW}}(X),G)
\cong
\prod_{\alpha\in I_n}G.
$$
The relative cohomology
$H^m(X^n,X^{n-1};G)$ is this product when
$m=n$ and is zero otherwise.

\noindent
The long exact sequences of the skeleta
give the following expanded version of the
preceding diagram.  The zeros express the
relative vanishing in the adjacent degrees.
{\small
$$
\xymatrix@C=.1em@R=1.3em{
  & & & 0 \ar[ld] \\
0
  & & H^{n-1}(X^n) \ar[ld] \\
  & H^{n-1}(X^{n-1})
    \ar[lu]
    \ar[ld]_{\partial_n^*} \\
H^n(X^n,X^{n-1})
  & & H^{n-1}(X^{n-1},X^{n-2})
      \ar[ll]_-{D^{n-1}}
      \ar[lu]_{j_{n-1}^*}
  & & H^{n-2}(X^{n-2},X^{n-3})
      \ar[ll]_-{D^{n-2}}
      \ar[ld]^{j_{n-2}^*} \\
  & & & H^{n-2}(X^{n-2})
          \ar[lu]^{\partial_{n-1}^*}
          \ar[ld] \\
  & & 0.
}
$$
}
The relative groups $H^k(X,X^{n+1};G)$
vanish for $k\leq n+1$.  Hence
$$
\xymatrix@C=.55em{
0 \ar[r]
  & H^n(X,X^{n+1};G) \ar[r]
  & H^n(X;G) \ar[r]^\cong
  & H^n(X^{n+1};G) \ar[r]
  & H^{n+1}(X,X^{n+1};G)
    \ar[r]
  & 0
}
$$
has zero outer relative terms.  Exactness
then identifies singular cohomology with
the kernel modulo image of the cellular
cochain maps in the expanded diagram.

\noindent
Finally, the evaluation maps from universal
coefficients give the commutative diagram
$$
\xymatrix@C=.55em{
H^n(X^n,X^{n-1};G)
  \ar[r] \ar[d]_h
  & H^n(X^n;G)
    \ar[r] \ar[d]_h
  & H^{n+1}(X^{n+1},X^n;G)
    \ar[d]_h \\
\operatorname{Hom}
(H_n(X^n,X^{n-1}),G)
  \ar[r]
  & \operatorname{Hom}(H_n(X^n),G)
    \ar[r]
  & \operatorname{Hom}
    (H_{n+1}(X^{n+1},X^n),G).
}
$$
The outer vertical maps are isomorphisms
because the relative homology groups are
free.  This identifies the cellular
cochain differential with the dual of the
cellular chain differential.
\end{proof}

\begin{theorem}[Mayer--Vietoris]
\label{at-theorem-mayer-vietoris-cohomology}
If $X=U\cup V$ is an excisive union, then
there is a natural long exact sequence
{\small
$$
\cdots\longrightarrow H^n(X;G)
\longrightarrow H^n(U;G)\oplus H^n(V;G)
\longrightarrow H^n(U\cap V;G)
\longrightarrow H^{n+1}(X;G)
\longrightarrow\cdots.
$$
}
\end{theorem}
\section{Cup products}
\label{at-section-cup-products}

\noindent
The product structure is what most visibly
distinguishes cohomology from homology.  In
this section the coefficients lie in a
commutative ring $R$ with unit.  We will
construct a product on
$$
\bigoplus_{n\geq0}H^n(X;R).
$$

\begin{definition}
\label{at-definition-cup-product}
Let $\varphi\in C^k(X;R)$ and
$\psi\in C^\ell(X;R)$.  Their {\it cup
product} is the cochain
$$
\varphi\smile\psi
\in C^{k+\ell}(X;R)
$$
defined on a singular simplex
$\sigma:\Delta^{k+\ell}\to X$ by
$$
\begin{aligned}
(\varphi\smile\psi)(\sigma)
&=
\varphi(\sigma|[v_0,\ldots,v_k])
\\
&\quad\cdot
\psi(\sigma|[v_k,\ldots,v_{k+\ell}]).
\end{aligned}
$$
Thus $\varphi$ reads the front face and
$\psi$ reads the back face, after which the
two values are multiplied in $R$.
\end{definition}

\noindent
The cup product is associative and
distributive directly from its definition.
The constant zero-cochain with value one is
a unit.  Hence
$$
\bigoplus_{n\geq0}C^n(X;R)
$$
is a graded ring under addition and cup
product.

\begin{remark}
\label{at-remark-graded-ring}
A {\it graded ring} is a direct sum
$$
A=\bigoplus_{i\geq0}A_i
$$
of abelian groups with
$$
A_iA_j\subset A_{i+j}.
$$
\end{remark}

\begin{proposition}[Leibniz rule]
\label{at-proposition-cup-leibniz}
If $\varphi$ has degree $k$, then
$$
\delta(\varphi\smile\psi)
=
\delta\varphi\smile\psi
+(-1)^k\varphi\smile\delta\psi.
$$
\end{proposition}

\begin{proof}
Evaluate both sides on a simplex and group
the faces before and after the common
vertex $v_k$.  All middle faces cancel in
pairs.  Moving the boundary past the first
$k$ vertices produces the sign.
\end{proof}

\noindent
It follows that a product of cocycles is a
cocycle.  A product of a coboundary and a
cocycle, in either order, is a coboundary.
Thus cup product is independent of the
chosen representatives and gives the
{\it cohomology ring}
$$
H^*(X;R)
=
\bigoplus_{n\geq0}H^n(X;R).
$$

\begin{proposition}
\label{at-proposition-cup-naturality}
For a continuous map $f:X\to Y$,
$$
f^*(\varphi\smile\psi)
=
f^*\varphi\smile f^*\psi.
$$
Consequently $f^*:H^*(Y;R)\to H^*(X;R)$
is a homomorphism of graded rings.
\end{proposition}

\begin{proposition}[Graded commutativity]
\label{at-proposition-cup-commutativity}
If $\alpha\in H^k(X;R)$ and
$\beta\in H^\ell(X;R)$, then
$$
\alpha\smile\beta
=
(-1)^{k\ell}\beta\smile\alpha.
$$
\end{proposition}

\noindent
In particular, the square of an odd-degree
class is killed by two.  It need not vanish
when the coefficient ring has two-torsion.
With coefficients in
$\mathbb{Z}/2\mathbb{Z}$ the sign
disappears.

\begin{example}[Real projective space]
\label{at-example-projective-mod-two-ring}
The cellular chain complex of
$\mathbb{R}P^n$ has one copy of
$\mathbb{Z}$ in each degree from zero to
$n$, with differentials alternating between
zero and multiplication by two.
$$
\xymatrix@C=.7em{
\mathbb{Z} \ar[r]^-{d_n}
  & \cdots \ar[r]
  & \mathbb{Z} \ar[r]^-2
  & \mathbb{Z} \ar[r]^-0
  & \mathbb{Z} \ar[r]
  & 0.
}
$$
After applying
$\operatorname{Hom}(-,
\mathbb{Z}/2\mathbb{Z})$, every
differential is zero.
$$
\xymatrix@C=.7em{
\mathbb{Z}/2\mathbb{Z}
  & \mathbb{Z}/2\mathbb{Z}
    \ar[l]_-0
  & \mathbb{Z}/2\mathbb{Z}
    \ar[l]_-0
  & \cdots \ar[l]
  & 0 \ar[l].
}
$$
Therefore there is one nonzero class in
every degree from zero through $n$.  If
$x$ is the degree-one class, then
$$
x^i
$$
is the nonzero class in degree $i$ and
$x^{n+1}=0$.  The product assertion is
proved geometrically below.  Thus
$$
H^*(\mathbb{R}P^n;
\mathbb{Z}/2\mathbb{Z})
\cong
(\mathbb{Z}/2\mathbb{Z})[x]/(x^{n+1}).
$$

\begin{proposition}
\label{at-proposition-cup-unions}
For a family of spaces,
$$
H^*\left(\coprod_iX_i;R\right)
\cong
\prod_iH^*(X_i;R).
$$
If the wedge point has a contractible
neighborhood, then
$$
\widetilde H^*\left(\bigvee_iX_i;R\right)
\cong
\prod_i\widetilde H^*(X_i;R).
$$
The reduced ring need not have a unit.
\end{proposition}

\begin{proof}
The additive statements were proved above.
Products of classes supported on distinct
components or distinct wedge summands
vanish, which identifies the ring
structures.
\end{proof}

\begin{example}
\label{at-example-wedge-projective-ring}
The spaces
$$
S^1\vee S^2
$$
and $\mathbb{R}P^2$ have isomorphic
mod-two cohomology groups.  Their rings are
not isomorphic.  If $a$ is the degree-one
class of the circle summand, then
$$
(a,0)^2=(a^2,0)=(0,0),
$$
whereas the square of the degree-one class
in $\mathbb{R}P^2$ is its nonzero
degree-two class.
\end{example}
\end{example}
\section{Tensor products and the Kunneth
theorem}
\label{at-section-kunneth}

\noindent
If $X$ and $Y$ are finite CW complexes,
then $X\times Y$ has the product CW
structure, whose $n$-cells are the products
$$
e^i\times e^j
$$
with $i+j=n$.  This motivates the tensor
product.

\begin{definition}
\label{at-definition-tensor-product}
Let $R$ be a commutative ring with unit and
let $M,N$ be $R$-modules.  A map
$M\times N\to A$ is {\it bilinear} if it
is $R$-linear in each variable.  The
{\it tensor product} $M\otimes_RN$ is an
$R$-module equipped with a bilinear map
having the universal property
$$
\xymatrix{
(m,n) \ar@{|->}[r]
  & m\otimes n \\
M\times N \ar[r]
  \ar[rd]_-b
  & M\otimes_RN
    \ar@{-->}[d]^{\overline b} \\
  & A.
}
$$
Explicitly, it is the free $R$-module on
symbols $m\otimes n$, modulo the relations
$$
\begin{aligned}
(m_1+m_2)\otimes n
&=m_1\otimes n+m_2\otimes n,
\\
m\otimes(n_1+n_2)
&=m\otimes n_1+m\otimes n_2,
\\
(rm)\otimes n
&=r(m\otimes n),
\\
m\otimes(rn)
&=r(m\otimes n).
\end{aligned}
$$
\end{definition}

\begin{proposition}
[Properties of tensor products]
\label{at-proposition-tensor-properties}
The following statements hold.
\begin{enumerate}
\item There is an isomorphism
$$
R\otimes_RM\cong M.
$$
\item There is a symmetry isomorphism
$$
M\otimes_RN\cong N\otimes_RM.
$$
\item Tensor product distributes over
arbitrary direct sums:
$$
\left(\bigoplus_{i\in I}M_i\right)
\otimes_RN
\cong
\bigoplus_{i\in I}(M_i\otimes_RN).
$$
\item A map $f:M_1\to M_2$ induces
$$
f\otimes1_N:
M_1\otimes_RN\longrightarrow M_2\otimes_RN.
$$
\item For fixed $N$, tensor product is a
covariant right-exact functor.
\item The tensor product of two free
$R$-modules is free.
\item Tensor products of graded modules
inherit a total grading.
\item Cellular chains satisfy
$$
C_*^{\mathrm{CW}}(X\times Y;R)
\cong
C_*^{\mathrm{CW}}(X;R)
\otimes_R C_*^{\mathrm{CW}}(Y;R).
$$
\end{enumerate}
\end{proposition}

\begin{proof}
The first four statements follow from the
universal property.  The fifth can be
pictured by the functor diagram
$$
\xymatrix{
  & \operatorname{Ab} \\
-\otimes_RN:
\operatorname{Mod}_R
  \ar[r]
  \ar[ur]
  & \operatorname{Mod}_R \ar[u].
}
$$
Indeed, a short exact sequence
$$
\xymatrix{
0 \ar[r]
  & M'' \ar[r]
  & M \ar[r]
  & M' \ar[r]
  & 0
}
$$
induces the right-exact sequence
$$
\xymatrix{
M''\otimes_RN \ar[r]
  & M\otimes_RN \ar[r]
  & M'\otimes_RN \ar[r]
  & 0.
}
$$
The left-derived functors are denoted
$\operatorname{Tor}^R_i(-,N)$.

\noindent
If
$$
L_1=\bigoplus_\alpha R
$$
and
$$
L_2=\bigoplus_\beta R,
$$
then
$$
L_1\otimes_RL_2
\cong
\bigoplus_{(\alpha,\beta)}R.
$$
For graded modules $A$ and $B$, define
$$
(A\otimes_RB)_n
=
\bigoplus_{i+j=n}A_i\otimes_RB_j.
$$
The summands lie on diagonals in the array
$$
\begin{matrix}
\vdots&&&&\\
A_0\otimes B_2&&&&\\
A_0\otimes B_1&A_1\otimes B_1&&&\\
A_0\otimes B_0&A_1\otimes B_0
&A_2\otimes B_0&\cdots&A_n\otimes B_0.
\end{matrix}
$$
The product CW structure identifies its
degree-$n$ basis with precisely these
pairs of cells.
\end{proof}

\noindent
Let $R$ be a principal ideal domain and let
$X$ and $Y$ be finite CW complexes.  The
cellular tensor product gives a natural
short exact sequence
{\small
$$
\xymatrix@C=.1em{
0 \ar[r]
  & \displaystyle\bigoplus_{i+j=k}
    H_i(X;R)\otimes_RH_j(Y;R)
    \ar[r]
  & H_k(X\times Y;R) \ar[r]
  & \displaystyle\bigoplus_{i+j=k-1}
    \operatorname{Tor}^R_1
    (H_i(X;R),H_j(Y;R))
    \ar[r]
  & 0.
}
$$
}

\noindent
The cohomological {\it cross product} is
defined by
$$
\begin{aligned}
H^*(X;R)\times H^*(Y;R)
&\longrightarrow H^*(X\times Y;R),
\\
(a,b)&\longmapsto
p_1^*(a)\smile p_2^*(b).
\end{aligned}
$$
It is bilinear and therefore factors
through the tensor product.
$$
\xymatrix{
H^*(X;R)\times H^*(Y;R)
  \ar[r]^-\times \ar[d]
  & H^*(X\times Y;R) \\
H^*(X;R)\otimes_RH^*(Y;R)
  \ar@{-->}[ur]_-\times.
}
$$
Give the tensor product the multiplication
$$
(a\otimes b)(c\otimes d)
=
(-1)^{|b||c|}(ac)\otimes(bd).
$$
Then the cross product is a homomorphism of
graded rings.

\begin{theorem}[Kunneth theorem]
\label{at-theorem-cohomology-kunneth}
Let $R$ be a field and let $X,Y$ be finite
CW complexes.  The cross product is an
isomorphism of graded rings:
$$
H^*(X;R)\otimes_RH^*(Y;R)
\xrightarrow{\cong}H^*(X\times Y;R).
$$
\end{theorem}

\noindent
The relative version gives a map
$$
H^*(X,A;R)\otimes_RH^*(Y,B;R)
\longrightarrow H^*(X\times Y,W;R),
$$
where
$$
W=(A\times Y)\cup(X\times B).
$$

\begin{proposition}
\label{at-proposition-projective-cup-ring}
For every $n\geq0$,
$$
H^*(\mathbb{R}P^n;
\mathbb{Z}/2\mathbb{Z})
\cong
(\mathbb{Z}/2\mathbb{Z})[x]/(x^{n+1}),
$$
where $|x|=1$.
\end{proposition}

\begin{proof}
Write $P^n=\mathbb{R}P^n$.  We already
know the additive groups.  It remains to
show that the product of nonzero classes in
degrees $i$ and $j$ is nonzero whenever
$i+j\leq n$.

\noindent
Use the disk model
$$
P^n=D^n/(x\sim-x)
$$
on the boundary.  For $i+j=n$, choose
projective subspaces $P^i,P^j\subset P^n$
whose coordinate blocks overlap in one
coordinate.  They meet in one point $p$.
Their open cells are transverse copies of
$\mathbb{R}^i$ and $\mathbb{R}^j$ inside
the open cell $\mathbb{R}^n$.

\noindent
Cup product with supports gives the
commutative comparison
$$
\xymatrix@C=.35em@R=2.2em{
H^i(P^n)\times H^j(P^n)
  \ar[r]^-\smile
  & H^n(P^n) \\
H^i(P^n,P^n-P^j)
  \times H^j(P^n,P^n-P^i)
  \ar[u]^\cong
  \ar[r]^-\smile
  \ar[d]_\cong
  & H^n(P^n,P^n-\{p\})
    \ar[u]_\cong
    \ar[d]^\cong \\
H^i(\mathbb{R}^n,
\mathbb{R}^n-\mathbb{R}^j)
  \times
H^j(\mathbb{R}^n,
\mathbb{R}^n-\mathbb{R}^i)
  \ar[r]^-\smile
  & H^n(\mathbb{R}^n,
    \mathbb{R}^n-\{0\}).
}
$$
All coefficients in this proof are
$\mathbb{Z}/2\mathbb{Z}$.  Excision gives
the downward isomorphisms.  Retraction onto
normal disks gives the diagram
$$
\xymatrix@C=.3em@R=2.2em{
H^i(\mathbb{R}^n,
\mathbb{R}^n-\mathbb{R}^j)
  \times
H^j(\mathbb{R}^n,
\mathbb{R}^n-\mathbb{R}^i)
  \ar[r]^-\smile
  \ar[d]_\cong
  & H^n(\mathbb{R}^n,
    \mathbb{R}^n-\{0\})
    \ar[d]^\cong \\
H^i(\mathbb{R}^i,
\mathbb{R}^i-\{0\})
  \times
H^j(\mathbb{R}^j,
\mathbb{R}^j-\{0\})
  \ar[r]_-\times
  & H^n(\mathbb{R}^n,
    \mathbb{R}^n-\{0\}).
}
$$
Each factor is one-dimensional.  The
Kunneth theorem says that the cross product
of its two generators is the generator on
the right.

\noindent
For completeness, the isomorphism on the
left can be tracked through
$$
\xymatrix@C=.35em@R=2em{
H^i(P^n)
  \ar[d]^\cong
  & H^i(P^n,P^{i-1})
    \ar[l]_\cong
    \ar[d]^\cong
  & H^i(P^n,P^n-P^j)
    \ar[l]_\cong
    \ar[r]^\cong
    \ar[d]^\cong
  & H^i(\mathbb{R}^n,
    \mathbb{R}^n-\mathbb{R}^j)
    \ar[d]^\cong \\
H^i(P^i)
  & H^i(P^i,P^{i-1})
    \ar[l]^\cong
  & H^i(P^i,P^i-\{p\})
    \ar[l]^\cong
    \ar[r]_\cong
  & H^i(\mathbb{R}^i,
    \mathbb{R}^i-\{0\}).
}
$$
The horizontal maps come from inclusions,
the rightmost maps from excision, and the
leftmost map from the cellular calculation.
This proves the product is nonzero when
$i+j=n$.  Naturality under
$P^{i+j}\hookrightarrow P^n$ gives the
general case.
\end{proof}

\begin{remark}
\label{at-remark-cup-intersection}
The preceding proof suggests that cup
product records intersections of suitably
represented submanifolds.
\end{remark}

\begin{remark}
\label{at-remark-map-to-projective-space}
For a map $f:X\to\mathbb{R}P^n$, the ring
homomorphism
$$
f^*:H^*(\mathbb{R}P^n;
\mathbb{Z}/2\mathbb{Z})
\longrightarrow H^*(X;
\mathbb{Z}/2\mathbb{Z})
$$
is determined by the image of $x$.
\end{remark}

\begin{remarks}
\label{at-remarks-projective-rings}
The same method gives the following rings.
\begin{enumerate}
\item
$$
H^*(\mathbb{R}P^\infty;
\mathbb{Z}/2\mathbb{Z})
\cong(\mathbb{Z}/2\mathbb{Z})[x],
$$
with $|x|=1$.
\item
$$
H^*(\mathbb{C}P^n;R)
\cong R[u]/(u^{n+1}),
$$
with $|u|=2$, and
$$
H^*(\mathbb{C}P^\infty;R)
\cong R[u].
$$
\item
$$
H^*(\mathbb{H}P^n;R)
\cong R[v]/(v^{n+1}),
$$
with $|v|=4$.
\item
$$
H^*(\mathbb{R}P^\infty
\times\mathbb{R}P^\infty;
\mathbb{Z}/2\mathbb{Z})
\cong
(\mathbb{Z}/2\mathbb{Z})[x,y].
$$
Finite products add one polynomial
generator for each factor.
\end{enumerate}
\end{remarks}

\begin{exercise}
Compute the integral cohomology ring of
$\mathbb{C}P^n$.
\end{exercise}

\begin{proof}
There is one cell in every even dimension
from $0$ through $2n$ and no odd cells.
Consequently its cellular chain complex is
$$
\xymatrix@C=.8em{
0\ar[r]&\mathbb{Z}\ar[r]
&0\ar[r]&\mathbb{Z}\ar[r]
&\cdots\ar[r]&\mathbb{Z}\ar[r]&0,
}
$$
with every differential zero.  This single
diagram replaces the two parity-dependent
chains in the exercise source, both of which
were false.

Let $\alpha\in H^2(\mathbb{C}P^n;\mathbb{Z})$
be the first Chern class of the dual
tautological line bundle.  Equivalently,
$\alpha$ is Poincare dual to a hyperplane.
Transverse hyperplanes meet in
$\mathbb{C}P^{n-k}$, so $\alpha^k$ is a
generator of $H^{2k}$ for $0\leq k\leq n$.
Dimensional reasons give $\alpha^{n+1}=0$.
Hence
$$
H^*(\mathbb{C}P^n;\mathbb{Z})
\cong
\mathbb{Z}[\alpha]/(\alpha^{n+1}),
\qquad |\alpha|=2.
$$
\end{proof}


\noindent
Integral universal coefficients give the
additive pattern
$$
\xymatrix@C=.65em{
H^*(P^{2m};\mathbb{Z})
  & \mathbb{Z} \ar[r]
  & 0 \ar[r]
  & \mathbb{Z}/2\mathbb{Z} \ar[r]
  & 0 \ar[r]
  & \cdots \ar[r]
  & \mathbb{Z}/2\mathbb{Z} \\
H^*(P^{2m+1};\mathbb{Z})
  & \mathbb{Z} \ar[r]
  & 0 \ar[r]
  & \mathbb{Z}/2\mathbb{Z} \ar[r]
  & 0 \ar[r]
  & \cdots \ar[r]
  & \mathbb{Z}.
}
$$
The entries run from degree zero on the
left to the top degree on the right.
Reduction modulo two gives
$$
\xymatrix@C=.55em@R=2em{
\mathbb{Z} \ar[r]
  \ar@{->>}[d]
  & 0 \ar[r] \ar[d]
  & \mathbb{Z}/2\mathbb{Z}
    \ar[r] \ar@{=}[d]
  & 0 \ar[r] \ar[d]
  & \cdots \ar[r]
  & H^n(P^n;\mathbb{Z}) \ar[d] \\
\mathbb{Z}/2\mathbb{Z} \ar[r]
  & \mathbb{Z}/2\mathbb{Z} \ar[r]
  & \mathbb{Z}/2\mathbb{Z} \ar[r]
  & \mathbb{Z}/2\mathbb{Z} \ar[r]
  & \cdots \ar[r]
  & \mathbb{Z}/2\mathbb{Z}.
}
$$
Since reduction respects cup products, the
degree-two integral class has the expected
nonzero powers.  More explicitly,
$$
H^*(P^{2m};\mathbb{Z})
\cong
\mathbb{Z}[a]/(2a,a^{m+1}),
$$
where $|a|=2$.  In odd dimension,
$$
H^*(P^{2m+1};\mathbb{Z})
\cong
\mathbb{Z}[a,b]/
(2a,a^{m+1},ab,b^2),
$$
where $|a|=2$ and $|b|=2m+1$.

\noindent
For $m>0$, the reduced integral cohomology
ring of
$$
P^{2m}\vee S^{2m+1}
$$
is the product of the reduced rings of the
two summands.  It is isomorphic to the
reduced integral cohomology ring of
$P^{2m+1}$.  With mod-two coefficients the
underlying groups still agree, but the
rings do not: the top power in projective
space is nonzero, while products between
wedge summands vanish.

\begin{theorem}[Hopf]
\label{at-theorem-hopf-division-dimension}
If $A$ is a finite-dimensional division
algebra over $\mathbb{R}$, then its real
dimension is a power of two.
\end{theorem}

\begin{proof}
Suppose $A\cong\mathbb{R}^n$.  Multiplying
on either side by a nonzero element is a
linear isomorphism.  Multiplication
therefore induces
$$
\varphi:P^{n-1}\times P^{n-1}
\longrightarrow P^{n-1}
$$
whose restriction to either factor is a
homeomorphism.  With mod-two coefficients,
$$
\varphi^*:
(\mathbb{Z}/2\mathbb{Z})[x]/(x^n)
\longrightarrow
(\mathbb{Z}/2\mathbb{Z})[x_1,x_2]/
(x_1^n,x_2^n).
$$
Restriction to the two factors shows that
$$
\varphi^*(x)=x_1+x_2.
$$
For example, the first restriction is the
commutative triangle
$$
\xymatrix{
H^*(P^{n-1}\times P^{n-1})
  \ar[d]_{i_1^*}
  & H^*(P^{n-1})
    \ar[l]_-{\varphi^*}
    \ar[ld]^\cong \\
H^*(P^{n-1}).
}
$$
It follows that
$$
0=\varphi^*(x^n)=(x_1+x_2)^n.
$$
Since $x_1^n=x_2^n=0$, every interior
binomial coefficient of order $n$ is even.
Write the binary expansion as
$$
n=2^{r_1}+\cdots+2^{r_s}.
$$
Over $\mathbb{Z}/2\mathbb{Z}$,
$$
(x_1+x_2)^n
=
\prod_{q=1}^s
(x_1^{2^{r_q}}+x_2^{2^{r_q}}).
$$
If $s>1$, selecting a nonempty proper
subset of factors produces a nonzero
interior monomial, a contradiction.  Hence
$s=1$ and $n$ is a power of two.
\end{proof}
\section{Poincare duality}
\label{section-poincare-duality}

\begin{definition}
\label{at-definition-topological-manifold}
An \emph{$n$-dimensional topological
manifold} is a Hausdorff,
second-countable space that is locally
homeomorphic to $\mathbb{R}^n$.  In this
section manifolds have no boundary.  Such a
manifold is \emph{closed} when it is
compact.
\end{definition}

\begin{theorem}[Poincare duality]
\label{at-theorem-poincare-integral}
Let $M$ be an oriented $n$-manifold.
Suppose that $M$ is connected and closed.
Then, for every $k$,
$$
H_k(M;\mathbb{Z})
\cong H^{n-k}(M;\mathbb{Z}).
$$
Every closed manifold, orientable or not,
satisfies
$$
H_k(M;\mathbb{Z}/2\mathbb{Z})
\cong
H^{n-k}(M;\mathbb{Z}/2\mathbb{Z}).
$$
\end{theorem}

\noindent
The precise orientation data will be
developed below.  Tensoring the integral
isomorphism with $\mathbb{Q}$, and then
using universal coefficients over a field,
gives
$$
\begin{aligned}
H_k(M;\mathbb{Q})
&\cong H_k(M;\mathbb{Z})
  \otimes\mathbb{Q} \\
&\cong H^{n-k}(M;\mathbb{Z})
  \otimes\mathbb{Q} \\
&\cong H^{n-k}(M;\mathbb{Q}) \\
&\cong H_{n-k}(M;\mathbb{Q}).
\end{aligned}
$$
Thus the free ranks are symmetric.  The
integral theorem and universal coefficients
also identify the torsion subgroup of
$H_k(M;\mathbb{Z})$ with that of
$H_{n-k-1}(M;\mathbb{Z})$.

\noindent
Let $x\in M$, and choose a coordinate
neighborhood $U$ of $x$.  Excision and the
long exact sequence of a pair give
$$
\begin{aligned}
H_n(M,M-\{x\};\mathbb{Z})
&\cong H_n(U,U-\{x\};\mathbb{Z}) \\
&\cong H_n(\mathbb{R}^n,
  \mathbb{R}^n-\{0\};\mathbb{Z}) \\
&\cong \widetilde H_{n-1}
  (\mathbb{R}^n-\{0\};\mathbb{Z}) \\
&\cong \widetilde H_{n-1}
  (S^{n-1};\mathbb{Z}) \\
&\cong\mathbb{Z}.
\end{aligned}
$$
This group is the local homology of $M$ at
$x$.

\begin{definition}
\label{at-definition-manifold-orientation}
An \emph{orientation of $\mathbb{R}^n$}
is a choice of generator of
$$
H_n(\mathbb{R}^n,
\mathbb{R}^n-\{0\};\mathbb{Z}).
$$
A \emph{local orientation} of $M$ at $x$
is a generator
$$
\mu_x\in
H_n(M,M-\{x\};\mathbb{Z}).
$$

An \emph{orientation of $M$} is a choice
of one local orientation $\mu_x$ at every
$x\in M$, locally compatible as follows.
For each sufficiently small coordinate ball
$B$, there must be one generator
$$
\mu_B\in H_n(M,M-B;\mathbb{Z})
$$
whose image at every $y\in B$ is $\mu_y$.
The relevant inclusion-induced maps are
$$
\xymatrix{
H_n(M,M-B;\mathbb{Z})
  \ar[r]^\cong \ar[d]_\cong
  & H_n(M,M-\{x\};\mathbb{Z}) \\
H_n(M,M-\{y\};\mathbb{Z}).
}
$$
The manifold is \emph{orientable} if it
admits an orientation.
\end{definition}

\begin{remark}
\label{at-remark-two-orientations}
A connected orientable manifold has exactly
two integral orientations.  Once the local
generator at one point is chosen,
compatibility determines all the others;
the two choices differ by sign.
\end{remark}

\begin{example}
\label{at-example-orientable-manifolds}
Every sphere $S^n$ is orientable.  The real
projective space $\mathbb{R}P^{2m}$ is not
orientable for $m\geq 1$, while
$\mathbb{R}P^{2m+1}$ is orientable.  The
Klein bottle is not orientable.
\end{example}

\begin{proposition}
\label{at-proposition-orientation-cover}
Every manifold $M$ has a canonical
two-sheeted covering
$$
p:\widetilde M\longrightarrow M
$$
whose total space is orientable.
\end{proposition}

\begin{proof}[Construction]
The fiber over $x$ consists of the two
generators of the local homology group:
$$
\widetilde M
=
\bigl\{(x,\mu)\mid
H_n(M,M-\{x\};\mathbb{Z})
=\mathbb{Z}\mu\bigr\}.
$$
The projection sends $(x,\mu)$ to $x$.
Over a small coordinate ball, continuation
of either local generator gives one sheet.
These sheets define the topology and show
that $p$ is a two-sheeted covering.  On the
total space, choose at $(x,\mu)$ the local
orientation represented by $\mu$.  This
choice varies compatibly, so
$\widetilde M$ is orientable.  For the
Mobius band, the construction unwinds the
twist and produces a cylinder.
\end{proof}

\begin{remark}
\label{at-remark-orientation-section}
An orientation of $M$ is precisely a
continuous section of its orientation
cover $p:\widetilde M\to M$.
\end{remark}

\begin{corollary}
\label{at-lemma-cover-components}
If $M$ is connected, then $M$ is
orientable if and only if $\widetilde M$
has two connected components.
\end{corollary}

\begin{proof}
An orientation gives a section, whose image
is one component; the opposite orientation
gives the other.  Conversely, if the
two-sheeted cover has two components, each
component maps homeomorphically to $M$ and
therefore supplies a section.
\end{proof}

\begin{corollary}
\label{at-lemma-no-index-two}
Let $M$ be connected.  If $\pi_1(M)$ has
no subgroup of index two, then $M$ is
orientable.
\end{corollary}

\begin{proof}
If $M$ were nonorientable, its orientation
cover would be connected.  Covering-space
theory would identify
$\pi_1(\widetilde M)$ with an index-two
subgroup of $\pi_1(M)$.
\end{proof}

\noindent
The orientation cover sits inside a larger
local-homology covering.  Its total space is
$$
M_{\mathbb{Z}}
=
\coprod_{x\in M}
H_n(M,M-\{x\};\mathbb{Z}),
$$
with the topology defined by continuation
over coordinate balls.  The projection
$q$ sends each local class over $x$ to
$x$.  The zero classes form a copy of $M$,
and the classes of absolute value $r$ form
a copy of $\widetilde M$ for each $r>0$.
Thus
$$
M_{\mathbb{Z}}
\cong M\sqcup
\coprod_{r\geq 1}\widetilde M.
$$
The maps fit into
$$
\xymatrix{
M_{\mathbb{Z}}\ar[dd]_q \\
&\widetilde M
  \ar@{^{(}->}[ul]
  \ar[dl]^p \\
M.
}
$$

\begin{definition}
\label{at-definition-ring-orientation}
Let $R$ be a commutative ring with unit.
An \emph{$R$-orientation} chooses at each
$x$ a generator of the rank-one
$R$-module
$$
H_n(M,M-\{x\};R),
$$
with the same local compatibility as above.
The manifold is \emph{$R$-orientable} if
such a choice exists.
\end{definition}

\noindent
The generators in a local fiber correspond
to the units $R^\times$.  When this group is
finite, their bundle is an
$|R^\times|$-sheeted covering, and an
$R$-orientation is a section.  In
$\mathbb{Z}/2\mathbb{Z}$ there is only one
unit, so every manifold is orientable with
these coefficients.  More generally, every
manifold is $R$-orientable when $2=0$ in
$R$.  The weaker condition that some
nonzero element is killed by two is not
enough: a local orientation must be a
generator, hence a unit.  An integral
orientation induces an $R$-orientation for
every coefficient ring $R$.

\begin{theorem}
\label{at-theorem-top-manifold-homology}
Let $M$ be a connected, closed
$n$-manifold.  For every $x\in M$, the map
$$
H_n(M;R)\longrightarrow
H_n(M,M-\{x\};R)
$$
is injective.  If $M$ is $R$-orientable,
it is an isomorphism.  If it is not
$R$-orientable, its image is
$$
\{r\in R\mid 2r=0\}
\subseteq H_n(M,M-\{x\};R)
\cong R.
$$
Moreover,
$$
H_j(M;R)=0
\qquad (j>n).
$$
\end{theorem}

\begin{proof}
Local classes continue along paths.  The
standard local-to-global argument for top
homology identifies the image of the map
with the classes fixed by continuation
around every loop, and also proves
injectivity.  The orientation monodromy of
a loop is either $1$ or $-1$.  In the
orientable case all loops act trivially.  In
the nonorientable case some loop acts by
$-1$, so its fixed classes are exactly
those satisfying $r=-r$.  Local homology
vanishes above degree $n$; the same
local-to-global argument, applied to a
finite coordinate-ball cover, gives the
last assertion.
\end{proof}

This is \cite[Theorem 3.26]{hat}.

\begin{lemma}
If $M$ is an $n$-manifold, then
$H_i(M;R)=0$ for $i>n$. This follows
componentwise from the closed and noncompact
cases.
\end{lemma}

\begin{proposition}
Let $M$ be a connected noncompact
$n$-manifold. Then
$$
H_i(M;R)=0
$$
for every $i\geq n$.
\end{proposition}

This is \cite[Proposition 3.29]{hat}.

\begin{remark}[Question]
\label{at-remark-local-determinant}
How does local homology relate to the line
$\bigwedge^n T_x^*M$ of a smooth manifold?
The integral local homology group is the
orientation line: choosing one of its two
generators chooses a sign.  A nonzero
element of $\bigwedge^n T_x^*M$ also
chooses a sign after one passes to its two
possible orientations.  These choices
correspond, but the abelian group
$H_n(M,M-\{x\};\mathbb{Z})$ is not
literally the real vector space
$\bigwedge^n T_x^*M$.
\end{remark}

\begin{definition}
\label{at-definition-fundamental-class}
Suppose $M$ is $R$-oriented.  A
\emph{fundamental class} is an element
$$
[M]\in H_n(M;R)
$$
whose image in every local group
$H_n(M,M-\{x\};R)$ is the chosen local
generator.
\end{definition}

\noindent
The preceding theorem gives a fundamental
class for every $R$-orientation.  Conversely,
such a class gives compatible local
generators and hence an orientation.  If
$B$ is a coordinate ball containing $x$,
the compatibility is visible in
$$
\xymatrix{
H_n(M;R)
  \ar[r]\ar[dr]
  &H_n(M,M-\{x\};R) \\
&H_n(M,M-B;R)\ar[u]^\cong.
}
$$

\begin{theorem}[Poincare duality]
\label{at-theorem-poincare-ring}
Let $M$ be a connected, closed,
$R$-oriented $n$-manifold.  Then
$$
H^k(M;R)\cong H_{n-k}(M;R)
$$
for every integer $k$.
\end{theorem}

\begin{proof}[Cellular picture]
For a triangulated manifold, a cell of
dimension $k$ has a dual cell of dimension
$n-k$.  With orientations chosen
compatibly, the incidence numbers for dual
cells are the transposes of the original
ones, up to the conventional signs.  Thus
the cochain complex of the original cell
structure identifies with the chain
complex of the dual structure, with degrees
reversed.  The source notes record the
degree reversal as
$$
\xymatrix@R=.8em{
C&C^* \\
0\ar[r]&n \\
1\ar[r]&n-1 \\
\vdots&\vdots \\
n\ar[r]&0.
}
$$
Passing to homology gives the displayed
isomorphism.  The cap-product formulation
below gives the intrinsic map and extends
beyond this cellular picture.
\end{proof}
\section{Cap products}
\label{section-cap-products}

\begin{definition}
\label{at-definition-cap-product}
Let $X$ be a space and $R$ a commutative
ring.  The \emph{cap product} is the
bilinear operation
$$
\frown:
C_k(X;R)\times C^\ell(X;R)
\longrightarrow C_{k-\ell}(X;R).
$$
For a singular simplex
$\sigma:[v_0,\ldots,v_k]\to X$, set
$$
\sigma\frown\varphi
=
\varphi
(\sigma|[v_0,\ldots,v_\ell])
\,\sigma|[v_\ell,\ldots,v_k],
$$
and extend $R$-bilinearly.
\end{definition}

\noindent
The boundary and coboundary satisfy
$$
\partial(\sigma\frown\varphi)
=(-1)^\ell
\bigl(
\partial\sigma\frown\varphi
-\sigma\frown\delta\varphi
\bigr).
$$
Consequently a cycle capped with a cocycle
is a cycle.  If either representative is
changed by a boundary or coboundary, the
result changes by a boundary.  Thus cap
product descends to
$$
\frown:
H_k(X;R)\times H^\ell(X;R)
\longrightarrow H_{k-\ell}(X;R).
$$

\begin{theorem}
\label{at-theorem-poincare-cap}
Let $M$ be a connected, closed,
$R$-oriented $n$-manifold.  Capping with
its fundamental class gives an isomorphism
$$
[M]\frown\mathord{-}:
H^k(M;R)\longrightarrow H_{n-k}(M;R)
$$
for every $k$.
\end{theorem}

\begin{proposition}
\label{at-proposition-cup-cap-adjunction}
For $\alpha\in C_{k+\ell}(X;R)$,
$\varphi\in C^k(X;R)$, and
$\psi\in C^\ell(X;R)$,
$$
\psi(\alpha\frown\varphi)
=(\varphi\smile\psi)(\alpha).
$$
\end{proposition}

\begin{proof}
Both sides evaluate $\varphi$ on the front
$k$-face, evaluate $\psi$ on the remaining
$\ell$-face, and multiply the two values.
\end{proof}

\noindent
In particular, the two chain maps
$$
\begin{aligned}
\varphi\smile\mathord{-}:
C^\ell(X;R)
&\longrightarrow C^{k+\ell}(X;R),\\
\mathord{-}\frown\varphi:
C_{k+\ell}(X;R)
&\longrightarrow C_\ell(X;R)
\end{aligned}
$$
are dual.  On homology, the universal
coefficient evaluation maps fit into the
commutative square
$$
\xymatrix@C=3.2em{
H^\ell(X;R)
  \ar[d]_{\varphi\smile\mathord{-}}
  \ar[r]^h
  &\operatorname{Hom}_R(H_\ell(X;R),R)
  \ar[d]^{
  (\mathord{-}\frown\varphi)^*} \\
H^{k+\ell}(X;R)
  \ar[r]^h
  &\operatorname{Hom}_R
  (H_{k+\ell}(X;R),R).
}
$$
Over a field the maps $h$ are
isomorphisms.  With integral coefficients
they are surjective, and $h$ in degree $r$
is an isomorphism when
$H_{r-1}(X;\mathbb{Z})$ is free abelian,
for example when it is finitely generated
and torsion-free.
\section{Consequences of duality}
\label{section-duality-consequences}

\begin{definition}
\label{at-definition-nonsingular-pairing}
Let $A$ and $B$ be $R$-modules.  A
bilinear pairing
$$
\lambda:A\times B\longrightarrow R
$$
is \emph{nonsingular} if both adjoint maps
$$
\begin{aligned}
A&\longrightarrow
\operatorname{Hom}_R(B,R),\\
a&\longmapsto
\bigl(b\longmapsto\lambda(a,b)\bigr),
\end{aligned}
$$
and
$$
\begin{aligned}
B&\longrightarrow
\operatorname{Hom}_R(A,R),\\
b&\longmapsto
\bigl(a\longmapsto\lambda(a,b)\bigr)
\end{aligned}
$$
are isomorphisms.
\end{definition}

\noindent
For a closed $R$-oriented $n$-manifold,
Poincare duality gives the intersection
pairing
$$
\begin{aligned}
H^k(M;R)\times H^{n-k}(M;R)
&\longrightarrow R,\\
(\varphi,\psi)
&\longmapsto
(\varphi\smile\psi)([M]).
\end{aligned}
$$
It is independent of the chosen cocycle
representatives.

\begin{proposition}
\label{at-proposition-duality-pairing}
The intersection pairing is nonsingular if
$R$ is a field.  For $R=\mathbb{Z}$, it
induces a nonsingular pairing after the
torsion subgroups of both cohomology groups
are factored out.
\end{proposition}

\begin{proof}
Let
$$
D=[M]\frown\mathord{-}:
H^k(M;R)\longrightarrow H_{n-k}(M;R).
$$
In the field case, both arrows in
$$
\xymatrix@C=2.6em{
H^{n-k}(M;R)\ar[r]^h
  &\operatorname{Hom}_R
  (H_{n-k}(M;R),R)
  \ar[r]^{D^*}
  &\operatorname{Hom}_R
  (H^k(M;R),R)
}
$$
are isomorphisms.  The composite sends
$\psi$ to the functional
$$
\varphi\longmapsto
\psi([M]\frown\varphi)
=(\varphi\smile\psi)([M]),
$$
so it is the second adjoint of the
intersection pairing.  The first adjoint
is handled in the same way, or by graded
commutativity.

For integral coefficients, the kernel of
the evaluation map $h$ is the torsion
subgroup supplied by the universal
coefficient theorem.  Poincare duality
preserves torsion.  The same argument
therefore becomes an isomorphism after
passing to torsion-free quotients.
\end{proof}

\begin{corollary}
\label{at-lemma-primitive-dual-class}
Assume the hypotheses of the proposition.
Let $\alpha$ be a class in $H^k(M;R)$.
It generates a rank-one direct summand if
and only if there is a class
$\beta\in H^{n-k}(M;R)$ for which
$$
(\alpha\smile\beta)([M])
$$
is a unit of $R$.  In the integral case,
the statement is understood in the
torsion-free quotient.
\end{corollary}

\begin{proof}
The class $\alpha$ generates a direct
summand precisely when some homomorphism to
$R$ takes $\alpha$ to a unit.  By
nonsingularity, every such homomorphism is
pairing with a suitable $\beta$.
\end{proof}

\begin{example}[Real projective space]
\label{at-example-rp-duality-product}
Put $P^n=\mathbb{R}P^n$ and use
$R=\mathbb{Z}/2\mathbb{Z}$.  The additive
cohomology groups are recorded by
$$
\xymatrix@C=.15em{
0\ar[r]&1\ar[r]&\ldots\ar[r]
  &i\ar[r]&\cdots\ar[r]
  &j\ar[r]&\cdots\ar[r]
  &i+j\ar[r]&\cdots\ar[r]&n\\
\mathbb{Z}/2\mathbb{Z}\ar[r]
  &\mathbb{Z}/2\mathbb{Z}\ar[r]
  &\ldots\ar[r]
  &\mathbb{Z}/2\mathbb{Z}\ar[r]
  &\cdots\ar[r]
  &\mathbb{Z}/2\mathbb{Z}\ar[r]
  &\cdots\ar[r]
  &\mathbb{Z}/2\mathbb{Z}\ar[r]
  &\cdots\ar[r]
  &\mathbb{Z}/2\mathbb{Z}.
}
$$
Let $x$ be the nonzero degree-one class.
In $P^{i+j}$, nonsingularity says that the
nonzero class $x^i$ pairs nontrivially with
the unique nonzero degree-$j$ class,
namely $x^j$.  Hence $x^{i+j}\neq 0$.
Naturality for
$P^{i+j}\hookrightarrow P^n$ gives the
same conclusion in $P^n$ whenever
$i+j\leq n$.  Thus duality recovers the
projective-space cup-product calculation.
\end{example}

\begin{example}[The torus]
\label{at-example-torus-pairing}
Let $T=S^1\times S^1$.  Its integral
cohomology has the pattern
$$
\xymatrix@R=.8em{
0&1&2\\
\mathbb{Z}\ar[r]
  &\mathbb{Z}\oplus\mathbb{Z}\ar[r]
  &\mathbb{Z}\\
&\left\langle{}\alpha\right\rangle{}
 \oplus
 \left\langle{}\beta\right\rangle{}
  &\left\langle{}\tau\right\rangle{}.
}
$$
Choose $\alpha$ to restrict to a generator
on a meridian, and choose $\beta$ from the
other circle factor.  The source records
this choice schematically as
$$
\xymatrix{
T&\mathbb{Z}\ar[r]
  &\mathbb{Z}\oplus\mathbb{Z}\ar[r]
  &\mathbb{Z}\\
S^1\ar@{^{(}->}[u]
  &\mathbb{Z}
  &\mathbb{Z}
   =\left\langle{}\alpha
    \right\rangle{}
  &0.
}
$$
Graded commutativity gives
$2\alpha^2=0$, and $H^2(T;\mathbb{Z})$
is torsion-free, so $\alpha^2=0$; this is
not a dimension argument.  If
$\gamma=n_1\alpha+n_2\beta$ satisfies
$\alpha\smile\gamma=\tau$, then
nonsingularity forces $n_2=\pm1$.
Reversing the choice of $\beta$ if needed,
we obtain
$$
\alpha\smile\beta=\tau,
\qquad
\beta\smile\alpha=-\tau.
$$
This agrees with the Kunneth description
of $H^*(S^1\times S^1;\mathbb{Z})$.
\end{example}
\section{Manifolds with boundary}
\label{section-manifolds-with-boundary}

\begin{definition}
\label{at-definition-boundary-manifold}
An \emph{$n$-manifold with boundary} is a
Hausdorff, second-countable space locally
homeomorphic to
$$
\mathbb{R}_+^n
=\{(x_1,\ldots,x_n)\in\mathbb{R}^n
\mid x_n\geq0\}.
$$
Its \emph{boundary} $\partial M$ consists
of the points corresponding locally to
$x_n=0$; its complement is the
\emph{interior}.
\end{definition}

\begin{remark}
\label{at-remark-oriented-boundary}
The boundary of an $n$-manifold is an
$(n-1)$-manifold without boundary.  If $M$
is oriented, its boundary inherits an
orientation.
\end{remark}

\begin{proposition}[Collar theorem]
\label{at-proposition-boundary-collar}
The boundary has a neighborhood in $M$
homeomorphic to
$$
\partial M\times[0,1),
$$
with $\partial M$ corresponding to
$\partial M\times\{0\}$.
\end{proposition}

\begin{definition}
\label{at-definition-boundary-orientation}
A manifold with boundary is
\emph{$R$-orientable} if its interior is
$R$-orientable.
\end{definition}

\begin{remark}[Question]
\label{at-remark-does-manifold-bound}
Given a closed manifold $M$, must there be
a compact manifold $W$ with
$\partial W=M$?  No: not every closed
manifold bounds.  This failure motivates
cobordism.
\end{remark}

\begin{definition}
\label{at-definition-unoriented-cobordism}
Two closed smooth $n$-manifolds $M$ and
$N$ are \emph{unoriented cobordant} if
there is a compact smooth
$(n+1)$-manifold $W$ such that
$$
\partial W\cong M\sqcup N.
$$
\end{definition}

\begin{remark}
\label{at-remark-cobordism-equivalence}
Unoriented cobordism is an equivalence
relation.  Cylinders give reflexivity,
relabeling the boundary gives symmetry,
and gluing cobordisms along their common
boundary gives transitivity.
\end{remark}

\begin{definition}
\label{at-definition-cobordism-ring}
The \emph{unoriented smooth cobordism ring}
is the graded set of cobordism classes
$$
\Omega_*^O
=\coprod_{n\geq0}\Omega_n^O,
$$
with addition induced by disjoint union and
multiplication induced by Cartesian product.
\end{definition}

\bibliography{my}
\bibliographystyle{amsalpha}

\end{document}
